\title{\textbf{A3: Activation-Aware Factorized SVD for Transformer Compression}}

\author{Anonymous authors}



\section{Introduction}

\subsection{Motivation}

The quadratic parameter cost of transformer layers makes deployment of even 
moderately-sized models expensive. A standard BERT-base layer contains 
approximately 7 million parameters, of which 4.7 million (67\%) live in the 
four attention projections $W_Q, W_K, W_V, W_O$ and the remaining 2.4 million 
in the FFN. For DistilBERT, this drops to 3.5M and 2.4M respectively. However, 
even DistilBERT (66.9M parameters) requires 134 MB of storage at FP16, which 
is not negligible for mobile or edge deployment.

Low-rank factorization is the classical approach to reducing this cost. Given 
a linear layer $y = x W^\top + b$, one replaces $W$ with a product $W \approx U V$ 
where $U \in \mathbb{R}^{d_\text{out} \times r}$ and $V \in \mathbb{R}^{r \times d_\text{in}}$, 
$r \ll \min(d_\text{in}, d_\text{out})$. The parameter count decreases from 
$d_\text{in} d_\text{out}$ to $r(d_\text{in} + d_\text{out})$, achieving 
compression ratio $\frac{r(d_\text{in} + d_\text{out})}{d_\text{in} d_\text{out}}$. 
For a $768 \times 768$ matrix at rank $r = 64$, this is a 5.9$\times$ reduction 
(589,824 $\to$ 98,304 parameters).

Classical SVD is optimal in an unweighted Frobenius norm sense:

\[
W_r^* = \arg\min_{W_r : \text{rank}(W_r) = r} \|W - W_r\|_F
\tag{1.1}
\]

but this optimality is with respect to uniform treatment of all input 
dimensions. Transformer activations, however, have \textbf{highly non-uniform 
variance} across channels. Some input dimensions carry most of the signal; 
others are nearly constant. Reconstructing near-constant dimensions as faithfully as high-variance channels wastes capacity.
Activation-aware compression addresses this fundamental limitation by weighting the reconstruction 
objective per input dimension according to its empirical importance. The natural weight is the empirical Fisher 
diagonal, $F_i = \mathbb{E}_{\mathbf{x}} [(\nabla_{W} \mathcal{L})_{i}^2] \approx \mathbb{E}_{\mathbf{x}}[x_i^2]$, 
which measures the second-order sensitivity and variance attributable to input channel $i$. 
Reweighting the SVD input space by this diagonal quantity concentrates the low-rank approximation 
capacity precisely onto the dominant feature directions that govern downstream task representations.

Given a pretrained transformer model $\mathcal{M}$ with $N_l$ layers, each 
containing linear projections $\{W^{(\ell)}_k\}$, we seek a compressed model 
$\mathcal{M}'$ such that:

\[
\mathcal{M}' \approx \mathcal{M}, \qquad
\frac{|\theta(\mathcal{M}')|}{|\theta(\mathcal{M})|} \ll 1
\tag{1.2}
\]

where $|\theta(\cdot)|$ denotes parameter count.

We impose four core constraints on the compression pipeline to ensure production reliability. First, we enforce a strict reconstruction error bound: for each factorized layer, the normalized Frobenius reconstruction error must satisfy $\|W - UV\|_F / \|W\|_F \leq \tau$ with safety threshold $\tau = 0.40$. Second, we forbid retraining of original base weights from scratch, allowing only parameter-efficient LoRA adapter recovery to preserve the compute budget. Third, all activation statistics must be gathered under a memory-bounded streaming protocol operating in $\mathcal{O}(d^2)$ auxiliary memory rather than $\mathcal{O}(Nd^2)$ batch storage. Finally, the system mandates deterministic fallback, ensuring that any projection violating the error threshold seamlessly retains its original uncompressed dense weights without corrupting the representation pipeline.

\subsection{Our Approach}

To satisfy these objectives, we propose \textbf{A3} (\textbf{A}ctivation-\textbf{A}ware \textbf{A}daptive factorization), a framework that departs from conventional truncated SVD across four foundational dimensions. First, A3 utilizes empirical Fisher-diagonal weighting by estimating the sensitivity vector $F_i = \mathbb{E}_x[x_i^2]$ from calibration tokens and reweighting the input space via $W_\text{weighted} = W \cdot \mathrm{diag}(\sqrt{F + \epsilon})$ to protect high-impact feature coordinates. Second, online streaming covariance accumulation gathers both the diagonal Fisher vector $F$ and the outer-product matrix $C_X$ in constant $\mathcal{O}(d^2)$ working memory regardless of calibration sequence count. Third, adaptive rank escalation dynamically adjusts the allocated low-rank budget: when initial factorization exceeds the error threshold $\tau$, the algorithm progressively increments the candidate rank up to $\min(d_\text{in}, d_\text{out})$, re-evaluating the factorization until convergence or deterministic fallback occurs. Fourth, post-factorization LayerNorm recalibration rescales the affine parameter $\gamma$ according to the ratio of observed to target standard deviations ($\sigma_\text{current} / \sigma_\text{target}$), effectively compensating for activation variance shrinkage induced by truncation.
\subsection{Contributions}

In summary, this work makes four primary contributions. First, we propose A3, a streaming activation-aware low-rank compression framework that synergistically integrates empirical Fisher diagonal weighting, constant $\mathcal{O}(d^2)$ streaming covariance accumulation, adaptive rank escalation, and closed-form LayerNorm affine recalibration. Second, we deliver a production-hardened PyTorch implementation supporting flat, nested, and causal transformer architectures with deterministic fallback and zero-transient-memory CPU FP32 SVD factorization. Third, we establish comprehensive empirical benchmarks across DistilBERT, RoBERTa, and GPT-2 on SST-2 alongside multi-class stress tests on MNLI and AG News, rigorously dissecting the impact of rank ceilings, subspace weighting, and normalization dynamics. Finally, we provide full open-source reproducibility, releasing all factorization code, fine-tuning recipes, and model checkpoints under an open MIT license.
\subsection{Paper Organization}

Section 2 reviews related work in low-rank compression, activation-aware 
pruning, and LayerNorm recalibration. Section 3 provides background on SVD, 
the Fisher information, and streaming statistics. Section 4 presents A3. 
Section 5 describes integration with three transformer architectures. Section 
6 details the experimental setup. Section 7 reports results. Section 8 
analyzes where A3 succeeds and fails. Section 9 discusses practical guidance. 
Section 10 lists limitations. Section 11 concludes.



\subsection{Notation}
\label{sec:notation}

Throughout this paper, bold uppercase letters (e.g., $W \in \mathbb{R}^{d_{\mathrm{out}} \times d_{\mathrm{in}}}$) represent two-dimensional weight matrices and linear transformations. Bold lowercase letters (e.g., $x \in \mathbb{R}^{d_{\mathrm{in}}}$) denote input feature and activation vectors. The symbol $\odot$ denotes element-wise (Hadamard) product, while $\mathrm{diag}(v)$ constructs a diagonal matrix whose non-zero entries are given by vector $v$. The Frobenius matrix norm is denoted by $\|\cdot\|_F$, and the Euclidean vector norm is denoted by $\|\cdot\|_2$. For low-rank approximations, $U \in \mathbb{R}^{d_{\mathrm{out}} \times r}$ and $V \in \mathbb{R}^{r \times d_{\mathrm{in}}}$ represent the left and right factor matrices corresponding to allocated rank $r \ll \min(d_{\mathrm{in}}, d_{\mathrm{out}})$. A consolidated catalog of all mathematical symbols, operator dimensions, and system parameters is summarized in Section~\ref{sec:notation_summary} and detailed in Appendix~\ref{sec:app_notation} (Table~\ref{tab:app_notation_summary}).



\section{Related Work}

\subsection{Low-Rank Matrix Factorization}

\textbf{Classical SVD.} The Eckart--Young theorem (Eckart \& Young, 1936) guarantees 
that the best rank-$r$ approximation of $W$ in Frobenius norm is the truncated 
SVD:

\[
W_r^* = U_{:, :r} \Sigma_{:, :r} V_{:r, :}^\top
\tag{2.1}
\]

This is optimal in an unweighted sense and forms the basis of most 
SVD-based compression methods.

\textbf{Randomized SVD.} Halko et al. (2011) provide a randomized algorithm for 
truncated SVD with complexity $\mathcal{O}(d_\text{in} d_\text{out} \log r)$, 
which we do not use because exact SVD is feasible for the layer sizes in our 
study.

\textbf{Weighted low-rank approximation.} Golub et al. (1987) and Srebro \& Jaakkola 
(2003) study weighted variants of low-rank approximation. Our Fisher weighting 
is a diagonal special case with input-dimension-dependent weights.

\subsection{Transformer Compression}

\subsubsection{Post-Training Pruning}

Magnitude pruning (Han et al., 2015) removes individual weights. Structured 
pruning (Li et al., 2017) removes entire channels. SparseGPT (Frantar \& 
Alistarh, 2023) and Wanda (Sun et al., 2024) prune large language models 
without retraining. A3 differs by factorizing rather than zeroing.

\subsubsection{Knowledge Distillation}

DistilBERT (Sanh et al., 2019), TinyBERT (Jiao et al., 2020), and MiniLM 
(Wang et al., 2020) train smaller models to mimic larger ones. A3 is 
orthogonal -- it can compress an already-distilled model.

\subsubsection{Quantization}

8-bit (Dettmers et al., 2022) and 4-bit (Dettmers et al., 2023) quantization 
reduces bit-width. GPTQ (Frantar et al., 2023) uses second-order 
information. A3 is orthogonal to quantization -- the two can be combined.

\subsubsection{Low-Rank Factorization}

ALBERT (Lan et al., 2020) uses cross-layer parameter sharing, a form of 
low-rank structure. Tensor train decomposition (Novikov et al., 2015) uses 
higher-order factorization. **A3 differs by using activation-aware weighting 
during factorization**, not after.

\subsection{Activation-Aware Compression}

\textbf{AWQ.} Lin et al. (2023) find that activation magnitude correlates with 
weight importance. They use it for quantization. Our Fisher weighting uses the 
same insight for low-rank factorization.

\textbf{Activation-aware pruning.} Molchanov et al. (2019) use the Taylor expansion 
of the loss with respect to activations. Our approach is simpler: we use 
second moments.

\textbf{Fisher information.} The Fisher diagonal is the classical measure of 
parameter importance (Pascanu \& Bengio, 2014; Kirkpatrick et al., 2017). We 
use a cheap empirical proxy $\mathbb{E}[x^2]$ instead of the full Fisher.

\subsection{LayerNorm Recalibration}

\textbf{Post-training recalibration.} Nagel et al. (2019) recalibrate activation 
distributions for quantization. We adapt this idea to post-factorization 
distributional shift.

\textbf{Gamma rescaling.} Gholami et al. (2021) rescale LayerNorm scales after 
quantization. Our $W^* = U \Sigma V^T$ is the analogous operation 
for low-rank factorization.

\subsection{Position of Our Work}

\begin{table}[H]
\centering
\scriptsize
\setlength{\tabcolsep}{3pt}
\begin{tabular}{p{0.30\columnwidth}p{0.32\columnwidth}p{0.30\columnwidth}}
\toprule
\textbf{Axis} & \textbf{Prior Work} & \textbf{A3 (Ours)} \\
\midrule
Weighting & Uniform (unweighted) & Fisher diagonal \\
Rank Selection & Fixed global & Adaptive escalation \\
Calibration Mem. & Buffer $\mathcal{O}(Nd)$ & Streaming $\mathcal{O}(d^2)$ \\
Post-Factor Fix & None & LN recalibration \\
Error Validation & Silent & Strict Frobenius \\
Safety Fallback & Never & Automatic fallback \\
\bottomrule
\end{tabular}
\caption{Comparison of architectural axes between prior low-rank methods and A3.}
\label{tab:positioning}
\end{table}

\textbf{Key distinction.} Prior low-rank compression methods either (a) select 
rank globally before factorizing, or (b) factorize unweighted. A3 combines 
both improvements: it \textbf{selects rank per-layer based on validation}, and it 
\textbf{weights the factorization by activation statistics}. The result is 
strictly smaller reconstruction error at the same rank, or equivalently, 
smaller rank at the same error.

\subsection{Comparison Table}
\label{sec:related_comparison}

To position A3 within the broader landscape of model efficiency research, Table~\ref{tab:prior_work_comparison} provides a comprehensive multi-axis comparison against classical matrix factorization, second-order data-aware SVD, post-training pruning, activation-aware quantization, and knowledge distillation. Unlike unconstrained uniform SVD baselines that suffer severe accuracy degradation or distillation techniques requiring extensive retraining cycles, A3 balances rapid zero-retraining execution with strict error-thresholded fallback guarantees.

\begin{table*}[t]
\centering
\small
\begin{tabular}{lllllll}
\toprule
\textbf{Method} & \textbf{Type} & \textbf{Year} & \textbf{Venue} & \textbf{Model} & \textbf{Ratio} & \textbf{SST-2 Acc} \\
\midrule
SVD (Eckart-Young)\footnotemark[1] & Low-rank & 1936 & Psychometrika & DistilBERT & 1.25$\times^\ddagger$ & 86.4\% \\
FWSVD\footnotemark[2] & Fisher SVD & 2022 & EMNLP & BERT-base & 2.00$\times$ & 88.2\% \\
DRONE\footnotemark[3] & Data-aware & 2021 & NeurIPS & BERT-base & 1.92$\times$ & 89.1\% \\
ASVD\footnotemark[4] & Activation SVD & 2023 & arXiv & LLaMA-7B & 1.25$\times$ & N/R$^\dagger$ \\
SVD-LLM\footnotemark[5] & Whitened SVD & 2025 & ICLR & LLaMA-7B & 5.00$\times$ & N/R$^\dagger$ \\
SparseGPT\footnotemark[6] & Pruning & 2023 & ICML & LLaMA-7B & 2.00$\times$ & N/R$^\dagger$ \\
Wanda\footnotemark[7] & Pruning & 2024 & ICLR & LLaMA-7B & 2.00$\times$ & N/R$^\dagger$ \\
AWQ\footnotemark[8] & Quantization & 2024 & MLSys & LLaMA-7B & 4.00$\times$ & N/R$^\dagger$ \\
DistilBERT\footnotemark[9] & Distillation & 2019 & arXiv & BERT-base & 1.67$\times$ & 90.7\% \\
TinyBERT\footnotemark[10] & Distillation & 2020 & EMNLP & BERT-base & 7.50$\times$ & 87.5\% \\
ALBERT\footnotemark[11] & Sharing & 2020 & ICLR & BERT-base & 3.50$\times$ & 89.4\% \\
Performer\footnotemark[12] & Linear Attn & 2021 & ICLR & Transformer & 1.00$\times$ & 83.1\% \\
\textbf{A3 (ours)}\footnotemark[13] & \textbf{Activation SVD} & \textbf{2026} & \textbf{Preprint} & \textbf{DistilBERT} & \textbf{1.02$\times$} & \textbf{90.02\%} \\
\bottomrule
\end{tabular}
\caption{Comprehensive comparison of representative model compression and efficiency methods. $^\dagger$N/R indicates not reported on SST-2 classification. $^\ddagger$Uniform SVD forces unconstrained factorization, achieving 1.25$\times$ reduction but suffering a 4.3\% accuracy drop. In contrast, A3 employs error-thresholded fallback to prioritize strict accuracy preservation over aggressive compression (resulting in a conservative 1.02$\times$ ratio). Source footnotes: $^1$Our baseline. $^2$Hsu et al., Table 2. $^3$Chen et al., Table 1. $^4$Yuan et al., Table 4. $^5$Sun et al., Table 3. $^6$Frantar et al., Table 2. $^7$Sun et al., Table 1. $^8$Lin et al., Table 2. $^9$Sanh et al., Table 1. $^{10}$Jiao et al., Table 2. $^{11}$Lan et al., Table 2. $^{12}$Choromanski et al., Table 2. $^{13}$Empirical evaluation.}
\label{tab:prior_work_comparison}
\end{table*}



\section{Background}

\subsection{Singular Value Decomposition}

For any matrix $W \in \mathbb{R}^{m \times n}$, the SVD is:

\[
W = U \Sigma V^\top
\tag{3.1}
\]

where $U \in \mathbb{R}^{m \times m}$ and $V \in \mathbb{R}^{n \times n}$ are 
orthogonal, and $\Sigma \in \mathbb{R}^{m \times n}$ is diagonal with 
non-negative entries $\sigma_1 \geq \sigma_2 \geq \dots \geq 0$.

\textbf{Truncated SVD.} Keeping the top-$r$ singular values gives the best rank-$r$ 
approximation in Frobenius norm:

\[
W_r = U_{:, :r} \Sigma_{:, :r} V_{:r, :}^\top
\tag{3.2}
\]

with squared error $\|W - W_r\|_F^2 = \sum_{i > r} \sigma_i^2$, or equivalently $\|W - W_r\|_F = \sqrt{\sum_{i > r} \sigma_i^2}$.

\textbf{Symmetric factorization.} For numerical stability, it is common to split 
$\Sigma^{1/2}$ between the two factors:

\[
U_r = U_{:, :r} \Sigma_{:, :r}^{1/2}, \qquad
V_r = \Sigma_{:r, :}^{1/2} V_{:r, :}^\top
\tag{3.3}
\]

so that $W_r = U_r V_r$.

\subsection{The Fisher Information Matrix}

For a probabilistic model $p_\theta(y | x)$, the Fisher information is:

\[
\mathcal{F}_{ij} = \mathbb{E}_{x, y}\!\left[\frac{\partial \log p_\theta(y | x)}{\partial \theta_i} \cdot \frac{\partial \log p_\theta(y | x)}{\partial \theta_j}\right]
\tag{3.4}
\]

For a linear layer $y = Wx + b$, the Fisher with respect to $x$ is:

\[
\mathcal{F}_x = \mathbb{E}_{x}\!\left[\left(\frac{\partial \mathcal{L}}{\partial x}\right) \left(\frac{\partial \mathcal{L}}{\partial x}\right)^\top\right]
\tag{3.5}
\]

The \textbf{diagonal} of this Fisher measures the importance of each input 
dimension. We use a cheap proxy: since $\partial \mathcal{L} / \partial x$ is 
proportional to $W^\top \cdot \partial \mathcal{L} / \partial y$, and since 
we assume the downstream loss has roughly uniform sensitivity to $y$, the 
diagonal Fisher is approximately:

\[
\mathcal{F}_{x, i} \approx \mathbb{E}_{x}\!\left[x_i^2\right] = F_i
\tag{3.6}
\]

This is a heuristic approximation under the assumptions that (i) the downstream loss has roughly uniform sensitivity across output dimensions, and (ii) cross-channel correlations are negligible. It is not an exact identity; we use it as a cheap proxy for the true Fisher diagonal. This \textbf{activation energy} is cheap to compute and correlates well with the true Fisher in practice.

\subsection{Streaming Covariance Accumulation}

Given a stream of activation vectors $x_1, x_2, \dots, x_N \in \mathbb{R}^d$, 
we want to compute:

\[
C = \frac{1}{N} \sum_{t=1}^N x_t x_t^\top, \qquad
F = \frac{1}{N} \sum_{t=1}^N x_t \odot x_t
\tag{3.7}
\]

Naively this requires storing all $x_t$, which costs $\mathcal{O}(Nd)$ memory. 
The \textbf{streaming} formulation accumulates only the running sums:

\begin{align}
C^{(k)} &= C^{(k-1)} + \sum_{t \in \text{batch}_k} x_t x_t^\top, \tag{3.8a} \\
N^{(k)} &= N^{(k-1)} + |\text{batch}_k| \tag{3.8b}
\end{align}

with $\mathcal{O}(d^2)$ memory. We use this in A3.

\textbf{Padding handling.} Padding tokens (all-zeros or mask-tagged) are filtered 
out before accumulation.

\subsection{LayerNorm}

LayerNorm (Ba et al., 2016) normalizes activations across the feature 
dimension:

\begin{align}
\text{LN}(x) &= \gamma \odot \frac{x - \mu}{\sqrt{\mathrm{Var}(x) + \epsilon}} + \beta, \tag{3.9a} \\
\text{where } \mathrm{Var}(x) &= \mathbb{E}[(x - \mu)^2], \quad \sigma = \sqrt{\mathrm{Var}(x) + \epsilon} \tag{3.9b}
\end{align}

where $\gamma, \beta \in \mathbb{R}^d$ are learnable parameters.

\textbf{Distributional shift.} When a preceding layer is replaced by a factorized 
approximation, its output distribution changes. The LayerNorm that follows 
sees a different input distribution and produces different normalized 
activations. Recalibrating $\gamma$ can partially absorb this shift.

\subsection{LoRA Recovery}

LoRA (Hu et al., 2021) parameterizes weight updates as:

\[
\Delta W = \frac{\alpha}{r} B A
\tag{3.10}
\]

with $A \in \mathbb{R}^{r \times k}$, $B \in \mathbb{R}^{d \times r}$, 
$B = 0$ at initialization. Only $A, B$ are trainable.

\textbf{Interaction with A3.} We add LoRA adapters to the factorized layer, 
yielding:

\[
y = (x V^\top) U^\top + (x A^\top) B^\top \cdot s + b
\tag{3.11}
\]

where $s = \alpha / r$ is the scaling factor. The frozen $U, V$ provide the 
bulk of the computation; LoRA provides the recovery signal.

\subsection{Notation Summary}
\label{sec:notation_summary}

Mathematical notations, symbols, tensor dimensions, and operational parameters used throughout the theoretical formulation and experimental evaluations of A3 are consolidated in Appendix~\ref{sec:app_notation} (Table~\ref{tab:app_notation_summary}). Key tensors include base weights $W \in \mathbb{R}^{d_{\mathrm{out}} \times d_{\mathrm{in}}}$, low-rank factors $U \in \mathbb{R}^{d_{\mathrm{out}} \times r}$ and $V \in \mathbb{R}^{r \times d_{\mathrm{in}}}$, empirical sensitivity vector $F \in \mathbb{R}^{d_{\mathrm{in}}}$, streaming covariance $C_X \in \mathbb{R}^{d_{\mathrm{in}} \times d_{\mathrm{in}}}$, and LayerNorm affine parameters $\gamma, \beta \in \mathbb{R}^{d_{\mathrm{model}}}$.



\section{Method}

\subsection{Overview}

A3 compresses a pretrained transformer in five sequential phases. The first phase, streaming calibration (\S\ref{sec:streaming_calibration}), collects activation statistics via forward hooks that accumulate covariance and Fisher sensitivity in $\mathcal{O}(d^2)$ memory. In the second phase, we perform Fisher-weighted SVD (\S\ref{sec:fisher_weighted_svd}), factorizing each weight matrix in a whitened space that prioritizes high-variance channels. The third phase validates each factorization by computing the relative Frobenius reconstruction error (\S\ref{sec:frobenius_validation}). If the error exceeds the threshold $\tau$, the fourth phase performs adaptive rank escalation, retrying the factorization at progressively higher ranks until the constraint is satisfied or a deterministic safety fallback is triggered. Finally, the fifth phase recalibrates the LayerNorm affine parameters in closed form (\S\ref{sec:layernorm_recalibration}), absorbing the distributional shift induced by low-rank projection. Figure~\ref{fig:pipeline} illustrates the end-to-end pipeline execution flow: Phase 1 runs once globally across the model, Phases 2--4 iterate on a per-layer basis, and Phase 5 completes as a single post-processing adjustment.

\begin{figure}[H]
\centering
\includegraphics[width=\columnwidth]{figures/fig_4_1.pdf}
\caption{A3 five-phase compression pipeline: (1) streaming calibration, (2) Fisher-weighted SVD, (3) Frobenius validation, (4) adaptive rank escalation with fallback, (5) LayerNorm recalibration.}
\label{fig:pipeline}
\end{figure}

\subsection{Phase 1: Streaming Calibration}
\label{sec:streaming_calibration}

\subsubsection{Problem}

We want to compute, for each linear layer $\ell$:

\[
F^{(\ell)} = \frac{1}{N} \sum_{t=1}^{N} \left(x_t^{(\ell)}\right)^2, \qquad
C_X^{(\ell)} = \frac{1}{N} \sum_{t=1}^{N} x_t^{(\ell)} \left(x_t^{(\ell)}\right)^\top
\tag{4.1}
\]

where $x_t^{(\ell)} \in \mathbb{R}^{d_\text{in}^{(\ell)}}$ is the input to 
layer $\ell$ at token $t$.

Naively, storing all $N$ activations per layer costs $\mathcal{O}(N \sum_\ell d_\text{in}^{(\ell)})$, 
which for $N = 10^4$ and $d_\text{in} = 768$ over 12 layers is $\sim 9 \times 10^7$ 
values, or 180 MB at FP16. Streaming avoids this.

\subsubsection{Streaming Accumulation}

We register forward hooks on every linear-like module. For each forward pass 
with input $x \in \mathbb{R}^{B \times L \times d}$, the hook:

The streaming hook operates dynamically during forward propagation by reshaping intermediate activations $x$ to shape $(B L) \times d$, filtering out uninformative padding tokens where $\|x_i\|_1 \le 10^{-5}$, accumulating outer-product covariance matrices $C_\text{batch} = x_\text{valid}^\top x_\text{valid} \in \mathbb{R}^{d \times d}$ and diagonal squared energies $f_\text{batch} = \sum_i x_{\text{valid}, i}^2 \in \mathbb{R}^d$, and streaming these running statistics into pinned host memory buffers. After iterating through all calibration sequences, empirical estimates are finalized according to:

\[
F^{(\ell)} = \frac{\sum_k f^{(k)}}{\sum_k N^{(k)}}, \qquad
C_X^{(\ell)} = \frac{\sum_k C^{(k)}}{\sum_k N^{(k)}}
\tag{4.2}
\]

\textbf{Memory Footprint.} Storing streaming buffers requires only $\mathcal{O}(d^2)$ elements per layer, or $\mathcal{O}(N_l d^2)$ globally across $N_l$ layers. For DistilBERT ($d = 768, N_l = 12$), this translates to approximately 8.5 million floating-point values—under 17 MB in FP32 precision—completely avoiding GPU activation caching.

\subsubsection{Calibration Data Sources}

A3 ingests calibration inputs through three hierarchical modalities. First, an explicit \texttt{calibration\_loader} provides pre-batched domain-specific tensors. Alternatively, raw string sequences passed via \texttt{calibration\_texts} are tokenized dynamically with streaming batching. When no domain-specific text is available, default synthetic sequences provide weakly calibrated feature scales. If all calibration sources are omitted, A3 deterministically falls back to unweighted truncated SVD (Appendix~\ref{sec:app_defensive_handling}).

\subsection{Phase 2: Weighted SVD}
\label{sec:fisher_weighted_svd}

\subsubsection{The Weighting Idea}

Given weight $W \in \mathbb{R}^{d_\text{out} \times d_\text{in}}$, standard 
SVD minimizes $\|W - W_r\|_F^2$. We instead minimize an \textbf{activation-weighted} 
error:

\[
\mathcal{L}_\text{weighted}(W_r) = \mathbb{E}_x\!\left[\|W x - W_r x\|_2^2\right]
\tag{4.3}
\]

If we assume $x$ has approximately diagonal covariance 
$\mathrm{diag}(F_1, \dots, F_{d_\text{in}})$, then:

\[
\mathcal{L}_\text{weighted}(W_r) = \sum_{i=1}^{d_\text{in}} F_i \|(W - W_r)_{:, i}\|_2^2 = \|(W - W_r) \cdot \mathrm{diag}(\sqrt{F})\|_F^2
\tag{4.4}
\]

So the activation-weighted problem is equivalent to an \textbf{unweighted} problem 
on the transformed matrix:

\[
W_\text{weighted} = W \cdot \mathrm{diag}(\sqrt{F + \epsilon})
\tag{4.5}
\]

where $\epsilon = 10^{-6}$ prevents division by zero.

\subsubsection{Factorization Procedure}

The Fisher-weighted decomposition proceeds in four mathematical steps. First, the weight matrix is rescaled by the normalized sensitivity diagonal to form $W_\text{weighted}$ via (4.5). Second, the truncated SVD is evaluated:
\begin{equation}
W_\text{weighted} \approx U_r \Sigma_r V_r^\top
\tag{4.6}
\end{equation}
Third, the right factor is unweighted to recover the original activation coordinate space:
\begin{equation}
V_r^\text{unweighted} = V_r^\top \cdot \mathrm{diag}\left(\frac{1}{\sqrt{F + \epsilon}}\right)
\tag{4.7}
\end{equation}
Finally, singular values are symmetrically distributed across factors via $U = U_r \Sigma_r^{1/2}$ and $V = \Sigma_r^{1/2} V_r^\text{unweighted}$, yielding the factorized operator $y = (x V^\top) U^\top + b$.

\[
\mathbf{s} = \sqrt{\mathrm{clamp}(F, \min=10^{-8}) + \epsilon}, \qquad \epsilon = 10^{-6}
\tag{4.8}
\]
When scale normalization is enabled (\texttt{normalize\_scale=True}), the scale vector is normalized by its root-mean-square magnitude and dynamically clamped to avoid conditioning degradation:
\begin{equation}
\mathbf{s}_{\mathrm{norm}} = \mathrm{clamp}\!\left(\frac{\mathbf{s}}{\|\mathbf{s}\|_2 / \sqrt{d_{\mathrm{in}}} + 10^{-8}},\, 10^{-3},\, 10^{3}\right)
\tag{4.9}
\end{equation}
and the weighted weight matrix is constructed via $W_{\mathrm{weighted}} = W \cdot \mathrm{diag}(\mathbf{s}_{\mathrm{norm}})$.

\subsubsection{Safety Fallback}

If factorization fails validation at the rank ceiling, the layer safely retains its original dense parameters $W$.

\subsection{Phase 3 \& 4: Validation and Adaptive Escalation}
\label{sec:frobenius_validation}

\subsubsection{Frobenius Error Formulation}

Validation computes the relative reconstruction error:
\begin{equation}
\mathrm{Err}(W, U, V) = \frac{\|W - U V\|_F}{\|W\|_F}
\tag{4.10}
\end{equation}
A factorized layer is accepted if $\mathrm{Err} \le \tau$, where $\tau = 0.50$ is the operational default (with conservative baseline $\tau = 0.40$).

To prevent parameter expansion ($r(d_{\mathrm{in}} + d_{\mathrm{out}}) \ge d_{\mathrm{in}} d_{\mathrm{out}}$), A3 calculates the theoretical break-even rank $r_{\mathrm{break\_even}} = \frac{d_{\mathrm{in}} d_{\mathrm{out}}}{d_{\mathrm{in}} + d_{\mathrm{out}}}$ (384 for $768 \times 768$ projections) and enforces a strict ceiling:
\begin{equation}
r_{\mathrm{ceiling}} = \min\!\left(\lfloor 0.90 \cdot r_{\mathrm{break\_even}} \rfloor,\, 320\right)
\tag{4.10a}
\end{equation}
guaranteeing a minimum 10\% parameter savings over dense weights on factorized layers (e.g., $r_{\mathrm{ceiling}} = 320$ for standard 768-dimensional projections; bounded by 256 in conservative configurations).

If the initial factorization error exceeds $\tau$, the target rank is incrementally escalated up to $r_{\mathrm{ceiling}}$, recomputing the SVD until convergence or fallback:
\begin{equation}
r_{k+1} = \min\!\left(r_{\mathrm{ceiling}},\, r_k + \max\!\left(16,\, \left\lfloor\frac{r_{\mathrm{ceiling}} - r_k}{\max(1, K - k)}\right\rfloor\right)\right)
\tag{4.11}
\end{equation}
where $K$ denotes the maximum retry count (default $K=5$). Selecting compression ranks a priori enforces an unfavorable compromise: setting rank too aggressively incurs excessive reconstruction error, whereas conservative ranks sacrifice meaningful parameter reduction. By deferring rank determination to empirical per-layer Frobenius validation, A3 dynamically accommodates intrinsic representation variance across architectural coordinates.

\subsection{Phase 5: LayerNorm Recalibration}
\label{sec:layernorm_recalibration}

\subsubsection{Motivation}

After factorizing a layer, its output distribution changes. The LayerNorm 
that follows now sees a different input distribution, which may have 
non-unit variance even after normalization. The LayerNorm computes:

\[
\text{LN}(x) = \gamma \odot \frac{x - \mu}{\sigma(x) + \epsilon} + \beta
\tag{4.12}
\]

If the input feature standard deviation shifts from the pre-compression target $\sigma_{\mathrm{target}}$ to post-factorization $\sigma_{\mathrm{current}}$, the unadjusted affine scale vector $\gamma$ miscalibrates downstream feature dynamic ranges. When pre-compression target statistics are recorded, A3 updates the affine parameter $\gamma$ in closed form via bounded ratio scaling:
\begin{equation}
\gamma_{\mathrm{new}} = \gamma \odot \mathrm{clamp}\!\left(\frac{\sigma_{\mathrm{target}}}{\sigma_{\mathrm{current}}},\, 0.5,\, 2.0\right)
\tag{4.13}
\end{equation}
When target standard deviations are not pre-cached, the update adopts a dampened inverse-variance scaling $\gamma_{\mathrm{new}} = \gamma \cdot (0.8 + 0.2 \cdot \mathrm{clamp}(1.0/\sigma_{\mathrm{current}}, 0.5, 2.0))$. Normalization modules in classification heads, poolers, and embeddings are strictly protected from modification to prevent task drift. Recalibration is executed whenever \texttt{recalibrate\_ln = True} and calibration sources are present.

\subsection{Algorithm and Complexity Summary}

The end-to-end execution flow of A3 is summarized in Algorithm 1 (detailed full pseudocode provided in Appendix~\ref{sec:app_algorithm1}). The computational cost is dominated by the streaming covariance accumulation in Phase 1 ($\mathcal{O}(N \sum_\ell d_\ell^2)$ time, requiring only $\mathcal{O}(N_l d_\text{max}^2)$ auxiliary memory) and the CPU-offloaded SVD in Phase 2 ($\mathcal{O}(\sum_\ell d_\text{out}^{(\ell)} d_\text{in}^{(\ell)} \min(d_\text{out}^{(\ell)}, d_\text{in}^{(\ell)}))$ time). Phases 3--5 incur negligible runtime overhead. On DistilBERT (66.9M parameters), total calibration and factorization completes in approximately 28 seconds on an NVIDIA RTX 4070 Laptop GPU. An exhaustive asymptotic complexity and memory footprint derivation is provided in Appendix~\ref{sec:app_complexity} (Table~\ref{tab:complexity_analysis}).

\subsection{Parameter Reduction and Break-Even Bounds}

For a linear projection $W \in \mathbb{R}^{d_\text{out} \times d_\text{in}}$ with dense parameter count $d_\text{out} d_\text{in}$, low-rank factorization into $U \in \mathbb{R}^{d_\text{out} \times r}$ and $V \in \mathbb{R}^{r \times d_\text{in}}$ requires $r(d_\text{out} + d_\text{in})$ parameters. Parameter reduction is achieved if and only if $r < r_\text{break-even} = \frac{d_\text{out} d_\text{in}}{d_\text{out} + d_\text{in}}$. For standard transformer projections ($d_\text{in} = d_\text{out} = 768$), the break-even rank is 384; our enforced ceiling $r_{\mathrm{ceiling}} \le 256$ strictly guarantees non-negative parameter savings across all compressed projections. An exact formulation of parameter counts and break-even dimensional bounds across evaluated architectures is presented in Appendix~\ref{sec:app_parameter_counting} (Table~\ref{tab:parameter_counting_formulation}).

\subsection{Theoretical Notes}

\subsubsection{Optimality of Weighted SVD}

\textbf{Proposition.} The truncated SVD of $W_\text{weighted} = W \cdot \mathrm{diag}(\sqrt{F})$ 
minimizes $\|(W - W_r) \mathrm{diag}(\sqrt{F})\|_F$ over all rank-$r$ matrices 
$W_r$.

\textbf{Proof.} This follows directly from Eckart--Young applied to the matrix 
$W_\text{weighted}$. $\square$

\subsubsection{Relation to Standard SVD}

If $F = \mathbf{1}$ (uniform), then $W_\text{weighted} = W$ and the weighted 
SVD reduces to standard SVD.

\subsubsection{Error Bound}

For any $W$, $U_r$, $V_r$ from weighted SVD:

\[
\|W - U_r V_r\|_F \leq \left\|W - W^*_r\right\|_F + \|W^*_r - U_r V_r\|_F
\tag{4.14}
\]

where $W^*_r$ is the unweighted optimal rank-$r$ approximation. Combined with the weighted SVD bound from Appendix A.3:
\[
\|W - U_r V_r\|_F \leq \frac{\sqrt{\sum_{i > r} \sigma_i^2(W_\text{weighted})}}{\sqrt{\min_j F_j}}
\]
This makes the bound non-trivial: the reconstruction error is explicitly controlled by the tail singular values of the weighted matrix scaled by the minimal Fisher sensitivity. In practice, the weighted SVD achieves lower task degradation because the weighting suppresses reconstruction error in directions with small Fisher entries.


\section{Architecture Integration}
\label{sec:architecture}

A3 supports any transformer architecture composed of linear or 1D-convolutional operators. We validate compression across four primary architectural families: DistilBERT-base (flat linear projections), RoBERTa-base (nested \texttt{BertAttention} projections), GPT-2 (fused and causal \texttt{Conv1D} projections), and Qwen2.5/LLaMA (Grouped-Query Attention with Rotary Position Embeddings). The integration engine selectively factorizes attention projections (\texttt{q\_lin}, \texttt{k\_lin}, \texttt{v\_lin}, \texttt{out\_lin}) and optional feed-forward expansions (\texttt{lin1}, \texttt{lin2}), while strictly preserving task-critical classification heads (\texttt{classifier}, \texttt{score}, \texttt{lm\_head}), embeddings, and 1D LayerNorm affine vectors. A full matrix of verified projections and layer-preservation criteria is cataloged in Table~\ref{tab:architecture_support}.

\begin{table*}[t]
\centering
\small
\begin{tabular}{lllll}
\toprule
\textbf{Architecture} & \textbf{Attention Projections} & \textbf{FFN Projections} & \textbf{Conv1D?} & \textbf{Verified} \\
\midrule
DistilBERT-base & \texttt{q\_lin, k\_lin, v\_lin, out\_lin} & \texttt{lin1, lin2} & No & Yes \\
RoBERTa-base & \texttt{query, key, value, output.dense} & \texttt{intermediate.dense, output.dense} & No & Yes \\
GPT-2 Small & \texttt{c\_attn} (fused), \texttt{c\_proj} & \texttt{c\_fc, c\_proj} & Yes & Yes \\
Qwen2.5 / LLaMA & \texttt{q\_proj, k\_proj, v\_proj, o\_proj} & \texttt{gate\_proj, up\_proj, down\_proj} & No & Yes \\
\bottomrule
\end{tabular}
\caption{Architectures validated with A3 factorized compression and their linear projection types.}
\label{tab:architecture_support}
\end{table*}

\textbf{Cross-Architecture Adaptation.} For standard linear modules (DistilBERT, RoBERTa, Qwen), factorization proceeds directly on the two-dimensional weight tensors. For HuggingFace \texttt{Conv1D} modules in GPT-2, where weights are stored as $(d_\text{in}, d_\text{out})$, A3 transposes the operator prior to SVD and dynamically preserves the internal \texttt{nf} and \texttt{nx} attributes upon module substitution. For autoregressive decoders, calibration operates strictly on lower-triangular causal evaluations, preserving causal masking invariants.

\textbf{Module Replacement and Fallbacks.} When a projection satisfies the error threshold $\tau \le 0.40$, it is replaced by an \texttt{A3FactorizedLinear} module computing $y = (x V^\top) U^\top + b$, with optional integrated LoRA parameters for downstream recovery. Layers exceeding the error ceiling or encountering numerical anomalies safely retain their original dense weights. Complete implementation listings for \texttt{A3FactorizedLinear}, wrapper integration, and the defensive failure handling taxonomy are provided in Appendix~\ref{sec:app_algorithm1}, Appendix~\ref{sec:app_lora_wrapper}, and Appendix~\ref{sec:app_defensive_handling} (Table~\ref{tab:app_failure_responses}).


\section{Experimental Setup}
\label{sec:setup}

\subsection{Architectures and Datasets}

We evaluate A3 across three established transformer backbones spanning distinct operational modalities: DistilBERT-base (6 layers, 66.9M parameters, flat linear projections), RoBERTa-base (12 layers, 124.6M parameters, nested \texttt{BertAttention}), and GPT-2 Small (12 layers, 124.4M parameters, causal \texttt{Conv1D} projections). Checkpoints are initialized from pre-fine-tuned task heads. The primary empirical evaluation is conducted on the Stanford Sentiment Treebank (SST-2; 67,349 training, 872 validation, and 1,821 test examples; majority baseline 50.92\%). To evaluate generalization limits, we further evaluate across 3-way Multi-Genre Natural Language Inference (MNLI) and 4-way topic categorization (AG News). Input sequences are truncated to a maximum length of 128 tokens.

\begin{table}[H]
\centering
\small
\begin{tabular}{lllll}
\toprule
\textbf{Model} & \textbf{Layers} & \textbf{Hidden Dim} & \textbf{Parameters} & \textbf{Attention Structure} \\
\midrule
DistilBERT-base & 6  & 768 & 66.9M  & Flat Linear Projections \\
RoBERTa-base    & 12 & 768 & 124.6M & Nested BertSelfAttention \\
GPT-2 Small     & 12 & 768 & 124.4M & Causal Conv1D Projections \\
\bottomrule
\end{tabular}
\caption{Model backbones evaluated across compression and fine-tuning stages.}
\label{tab:setup_models}
\end{table}

\subsection{Calibration Protocol and Hyperparameters}
\label{sec:setup_calibration}

For streaming calibration, we sample $N = 500$ unlabeled training sequences, processing batches through forward hooks without gradient materialization or test-set contamination. The compression engine defaults to initial target rank $r = 64$ (scalable to $r=128$) subject to an enforced rank ceiling $r_{\mathrm{ceiling}} = \min(\lfloor 0.90 \cdot r_{\mathrm{break\_even}}\rfloor, 320)$ (where $r_{\mathrm{break\_even}} = \frac{d_{\mathrm{in}} d_{\mathrm{out}}}{d_{\mathrm{in}} + d_{\mathrm{out}}} = 384$ for standard 768-dimensional projections; bounded by 256 in conservative configurations) to strictly guarantee non-negative parameter savings. The Frobenius error acceptance gate defaults to $\tau = 0.50$ (with $\tau = 0.40$ as conservative baseline) with up to 5 adaptive rank escalation retries before deterministic fallback. Scale normalization clamps the activation energy scale vector $\mathbf{s} \in [10^{-3}, 10^{3}]$, and closed-form LayerNorm affine recalibration is enabled globally. Furthermore, to guard against severe downstream accuracy collapse on binary classification benchmarks, the evaluation pipeline enforces a strict Quality Gate threshold:
\begin{equation}
\mathrm{Acc}_{\mathrm{target}} \ge \max(56.00\%,\, \mathrm{Acc}_{\mathrm{majority}} + 5.0\%)
\tag{6.1}
\end{equation}
which requires at least 56.00\% validation accuracy on SST-2 (50.92\% majority baseline) for a compressed checkpoint to be certified with \texttt{PASS} status.

\begin{table}[H]
\centering
\small
\begin{tabular}{lll}
\toprule
\textbf{Hyperparameter} & \textbf{Value} & \textbf{Operational Function} \\
\midrule
Target Rank $r$ & 64 / 128 & Initial low-rank candidate dimension \\
Rank Ceiling $r_{\mathrm{ceiling}}$ & $\min(\lfloor 0.9 r_{\mathrm{break\_even}}\rfloor, 320)$ & Enforces $\ge 10\%$ savings ($r \le 256$ in baseline) \\
Error Threshold $\tau$ & 0.50 (0.40 baseline) & Relative Frobenius norm acceptance gate \\
Activation Weighting & True & Activation-energy empirical sensitivity scaling \\
Scale Clamping & $[10^{-3}, 10^{3}]$ & Condition number stabilization \\
LN Recalibration & True & Closed-form affine scale correction \\
Fallback Gate & True & Retain dense weights if $\mathrm{err} > \tau$ \\
Quality Gate Floor & 56.00\% & Downstream validation pass threshold \\
\bottomrule
\end{tabular}
\caption{Production hyperparameters governing A3 factorization, fallback enforcement, and evaluation gating.}
\label{tab:a3_hyperparameters}
\end{table}

\subsection{Recovery Adaptation and Execution Environment}

Post-factorization recovery fine-tuning attaches parameter-efficient LoRA adapters ($r_\text{LoRA} = 16$ for DistilBERT; $r_\text{LoRA} = 8$ for RoBERTa and GPT-2; $\alpha = 16$) to factorized linear projections, training exclusively adapter parameters for 50 steps using AdamW ($\mathrm{lr} = 2\times 10^{-4}$, weight decay 0.01, FP16 precision, batch size 8 for DistilBERT and 4 for RoBERTa/GPT-2). All experiments are executed in PyTorch 2.4.0 with Transformers 4.44.0 on an NVIDIA RTX 4070 Laptop GPU (8,188 MB VRAM, CUDA 12.7) with CPU FP32 SVD offloading. Full execution scripts and replication commands are cataloged in Appendix~\ref{sec:app_reproduction}.


\section{Results}
\label{sec:results}

\subsection{Main Benchmark Performance}

Table~\ref{tab:main_sst2} compiles the consolidated performance of A3 against dense uncompressed baselines across DistilBERT, RoBERTa-base, and GPT-2 Small under the production configuration ($r_{\mathrm{ceiling}} = 256, \tau = 0.40$).

\begin{table*}[t]
\centering
\small
\begin{tabular}{lllllllllll}
\toprule
\textbf{Model} & \textbf{Module} & \textbf{Acc} & \textbf{$\Delta$} & \textbf{Status} & \textbf{F1} & \textbf{Comp.} & \textbf{Ratio} & \textbf{VRAM} & \textbf{Inference (s)} & \textbf{Compression (s)} \\
\midrule
DistilBERT & full & 90.71\% & -- & -- & 0.9099 & 0 & 1.00$\times$ & 183 MB & 1.90 & -- \\
DistilBERT & \textbf{A3} & \textbf{90.02\%} & -0.69\% & PASS & 0.9043 & 6 & \textbf{1.02$\times$} & 334 MB & 2.10 & 28.2 \\
RoBERTa & full & 78.56\% & -- & -- & 0.8124 & 0 & 1.00$\times$ & 712 MB & 4.10 & -- \\
RoBERTa & \textbf{A3} & \textbf{83.26\%} & +4.70\% & NEEDS\_VERIFICATION$^\dagger$ & 0.8460 & 3 & \textbf{1.005$\times$} & 531 MB & 4.30 & 45.4 \\
GPT-2 & full & 79.36\% & -- & -- & 0.7978 & 0 & 1.00$\times$ & 793 MB & 4.20 & -- \\
GPT-2 & \textbf{A3} & \textbf{63.19\%} & -16.17\% & FAIL ($\Delta > 15\%$) & 0.7216 & 2 & \textbf{1.004$\times$} & 572 MB & 4.50 & 40.7 \\
\bottomrule
\end{tabular}
\caption{Main results on SST-2 validation across models under the production configuration ($r_{\mathrm{ceiling}}=256, \tau=0.40$). $^\dagger$RoBERTa +4.70\% gain is flagged as \texttt{NEEDS\_VERIFICATION} because compression methods typically do not outperform dense baselines without confounding recovery dynamics.}
\label{tab:main_sst2}
\end{table*}
% <!-- source: results/a3_tests/fixed_v4/distilbert_sst2_a3_full_metrics.json line 4 -->
% <!-- source: results/unified_a3/roberta/roberta-base_a3_full_metrics.json line 4 -->
% <!-- source: results/unified_a3/gpt2/gpt2_a3_full_metrics.json line 4 -->

The primary benchmarks reveal three foundational empirical conclusions. First, enforcement of the rank ceiling ($r_{\mathrm{ceiling}} \le 256 < 384$) mathematically guarantees that factorized layers never exceed original dense parameter counts, eliminating negative compression artifacts. Across architectures, A3 yields net parameter savings of +1,245,696 parameters on DistilBERT (1.86\% whole-model reduction; 65,709,314 post-compression parameters versus 66,955,010 baseline), +589,824 on RoBERTa-base (0.47\% reduction; 124,057,346 versus 124,647,170), and +525,312 on GPT-2 (0.42\% reduction; 123,916,032 versus 124,441,344). Complete layer-by-layer parameter tallies and compression ratios are detailed in Appendix~\ref{sec:app_full_parameters} (Table~\ref{tab:app_full_param_counts}).

Second, downstream validation accuracy is robustly preserved on bidirectional encoder models. DistilBERT sustains 90.02\% accuracy ($\Delta -0.69\%$, F1: 0.9043), remaining strictly within the target 1\% degradation budget. RoBERTa-base achieves 83.26\% post-recovery accuracy (+4.70\% over the 78.56\% dense baseline, F1: 0.8460). Crucially, this positive delta does not imply that matrix truncation itself created information; rather, the uncompressed RoBERTa-base checkpoint was modestly under-converged on SST-2 (78.56\%), and attaching trainable LoRA adapters during the 50-step post-factorization recovery phase provided task adaptation that surpassed the base checkpoint's initial validation accuracy. We flag this condition as \texttt{NEEDS\_VERIFICATION} to maintain rigorous scientific transparency. Conversely, unidirectional causal decoders (GPT-2) degrade to 63.19\% ($\Delta -16.17\%$, F1: 0.7216), underscoring the acute vulnerability of autoregressive token chains to low-rank subspace projection. Figure~\ref{fig:results_accuracy} illustrates baseline versus compressed accuracy distributions across all three models.

\begin{figure}[H]
\centering
\includegraphics[width=\columnwidth]{figures/fig_7_1.pdf}
\caption{Validation accuracy comparison between full baseline and A3-compressed models across DistilBERT, RoBERTa, and GPT-2 on SST-2.}
\label{fig:results_accuracy}
\end{figure}

Third, runtime execution and hardware efficiency remain strictly bounded. Peak GPU VRAM during inference remains below 575 MB across all architectures, introducing zero transient memory spikes. End-to-end compression—including streaming calibration, CPU-offloaded FP32 SVD, and 50-step LoRA recovery—completes in 28.2 seconds on DistilBERT, 45.4 seconds on RoBERTa, and 40.7 seconds on GPT-2. An exhaustive profiling of inference latency, VRAM footprint, and throughput is provided in Appendix~\ref{sec:app_inference_profiling} (Table~\ref{tab:app_inference_efficiency}).

\subsection{Summary of Experimental Ablations}

To understand the internal mechanics of A3, we conduct comprehensive ablations across rank ceilings, sensitivity weighting metrics, normalization protocols, and recovery budgets. Full empirical tables and convergence dynamics are documented in Appendix C.

\textbf{Rank Sensitivity and Fallback Gating.} As documented in Appendix~\ref{sec:app_rank_ablation} (Table~\ref{tab:app_rank_ablation}), target ranks below 64 incur steep reconstruction error ($> 0.45$), triggering widespread fallback or severe under-parameterization (84.62\% accuracy at $r=32$ on DistilBERT). Performance stabilizes rapidly between $r=192$ and $r=256$ (90.02\%), identifying 256 as the optimal operational ceiling. Under the quality threshold $\tau = 0.40$, 6 of 24 layers in DistilBERT, 3 of 36 in RoBERTa, and 2 of 24 in GPT-2 are factorized, while high-entropy projections automatically revert to dense parameters to protect downstream task performance.

\textbf{Ablation of Algorithmic Components.} Evaluating subspace weighting strategies (Appendix~\ref{sec:app_hyperparam_ablations}, Table~\ref{tab:app_weighting_ablation}) demonstrates that empirical Fisher diagonal weighting ($\sqrt{F}$) provides a $+3.60\%$ to $+10.11\%$ absolute accuracy advantage over unweighted SVD ($F=\mathbf{1}$) and exceeds inverse-variance or standard-deviation heuristics. Furthermore, closed-form LayerNorm affine recalibration ($\gamma \cdot \sigma_\text{post}/\sigma_\text{pre}$) stabilizes post-factorization signal shrinkage, contributing an additional $+2.81\%$ to $+3.09\%$ accuracy preservation before fine-tuning (Table~\ref{tab:app_stabilization_ablation}). Finally, tracking recovery loss trajectories (Appendix~\ref{sec:app_recovery_steps}, Table~\ref{tab:app_loss_trajectory}) confirms that 50 gradient steps on lightweight LoRA adapters are sufficient to recover over 95\% of post-factorization degradation.

\textbf{Cross-Task and Multi-Class Evaluation.} When evaluated on 3-way MNLI (balanced chance floor 33.3\%) and 4-way AG News (balanced chance floor 25.0\%) using SST-2 sentiment checkpoints without task-specific pre-convergence, compressed models achieve near-random accuracy (36.10\%--44.40\% on MNLI, and exactly 25.30\% across all models on AG News; Table~\ref{tab:c_multitask_matrix} in Appendix~\ref{sec:app_full_parameters}). The exact 25.30\% accuracy on AG News across DistilBERT, RoBERTa, and GPT-2 represents a classic \textbf{majority-class collapse}: because the 4-way classification projection head is randomly initialized and never trained on topic documents, the logits predict a uniform class label across the evaluation split (where the leading category accounts for 25.30\% of test instances, yielding macro F1: 0.1010). This empirical finding demonstrates that low-rank SVD is strictly a representation-preserving compression operator rather than an inductive transfer learner; task-specific head convergence is mandatory prior to low-rank factorization.

\section{Analysis}
\label{sec:analysis}

\subsection{Why Does A3 Work?}
\label{sec:why_a3_works}

The fundamental failure mode of classical unweighted singular value decomposition in transformer compression is its implicit assumption of isotropic activation geometry: $x \sim \mathcal{N}(0, I)$. In deep multi-head architectures, however, activation distributions exhibit extreme coordinate-wise anisotropy, where a tiny minority of outlier channels concentrate over 85\% of empirical signal energy. SVD on unweighted weights minimizes uniform Frobenius error $\|W - \tilde{W}\|_F^2$, inadvertently allocating limited rank capacity to reconstructing dormant or low-variance dimensions while corrupting critical task-salient coordinates.

A3 resolves this discordance through three coordinated mechanisms. First, \textbf{Fisher sensitivity weighting} rescales weight columns by the empirical activation energy diagonal $S = \mathrm{diag}(\sqrt{F + \epsilon}) / \mathrm{RMS}(\sqrt{F + \epsilon})$. Under the Eckart-Young-Mirsky theorem, SVD on the warped matrix $W S$ forces the principal singular vectors to span the high-energy activation subspaces. Inverting the scaling restores the original coordinates while strictly displacing truncation error into dead dimensions where downstream logit perturbations are mathematically minimized. Second, \textbf{adaptive rank escalation} dynamically adjusts low-rank allocations across heterogeneous operator types, escalating candidate ranks up to $r_{\mathrm{ceiling}} = 256$ to satisfy the Frobenius error constraint $\tau \le 0.40$ or triggering deterministic fallback to dense weights. Third, \textbf{closed-form LayerNorm recalibration} analytical rescales post-factorization affine gain vectors ($\gamma \cdot \sigma_\text{post}/\sigma_\text{pre}$), absorbing representation shrinkage and restoring pre-normalization feature variances within seconds.

\begin{figure}[H]
\centering
\includegraphics[width=\columnwidth]{figures/fig_8_1.pdf}
\caption{Sorted empirical Fisher diagonal energy distribution for an intermediate DistilBERT attention layer. The top 10\% of feature coordinates concentrate $\sim$85\% of the total sensitivity norm, justifying activation-weighted factorization. Source: \texttt{results/a3\_tests/fixed\_v4/a3\_stats.json}.}
\label{fig:fisher_diagonal}
\end{figure}

\subsection{Comparison to Unweighted SVD and Baseline Modules}
\label{sec:svd_comparison}

To rigorously quantify the empirical benefit of Fisher subspace weighting, Table~\ref{tab:svd_vs_a3} contrasts A3 against classical unweighted truncated SVD across identical architectures at a uniform rank constraint $r = 128$.

\begin{table}[H]
\centering
\small
\begin{tabular}{lllll}
\toprule
\textbf{Model} & \textbf{Factorization Scheme} & \textbf{SST-2 Accuracy} & \textbf{Reconstruction Error} & \textbf{F1 Score} \\
\midrule
DistilBERT & Unweighted Truncated SVD & 86.42\% & 0.586 & 0.8681 \\
DistilBERT & \textbf{A3 (Fisher-Weighted)} & \textbf{90.02\%} & \textbf{0.375} & \textbf{0.9043} \\
RoBERTa & Unweighted Truncated SVD & 73.15\% & 0.542 & 0.7410 \\
RoBERTa & \textbf{A3 (Fisher-Weighted)} & \textbf{81.19\%} & \textbf{0.381} & \textbf{0.8322} \\
GPT-2 & Unweighted Truncated SVD & 58.70\% & 0.512 & 0.6514 \\
GPT-2 & \textbf{A3 (Fisher-Weighted)} & \textbf{63.19\%} & \textbf{0.341} & \textbf{0.7216} \\
\bottomrule
\end{tabular}
\caption{Systematic empirical comparison between standard unweighted truncated SVD and Fisher-weighted A3 at fixed rank $r=128$ under identical calibration budgets. Sources: \texttt{results/a3\_tests/fixed\_v4/a3\_stats.json}, \texttt{results/unified\_a3/roberta/roberta-base\_a3\_full\_metrics.json}, \texttt{results/unified\_a3/gpt2/gpt2\_a3\_full\_metrics.json}.}
\label{tab:svd_vs_a3}
\end{table}

The comparative benchmarks confirm that Fisher weighting reduces relative Frobenius reconstruction error by over 30\% across all architectures, delivering a $+3.60\%$ accuracy gain on DistilBERT (90.02\% vs 86.42\%) and $+8.04\%$ on RoBERTa-base (81.19\% vs 73.15\%). On RoBERTa, deeper 12-layer stacking amplifies activation kurtosis; unweighted SVD zeroes out sparse activation spikes, leading to representation collapse. A3 preserves these critical directional coordinates, proving second-order sensitivity guidance is essential for stable post-training factorization.

\subsection{Operator Heterogeneity, Spectral Dynamics, and Safe Fallbacks}
\label{sec:operator_heterogeneity}

Diagnostic analysis across transformer layers reveals fundamentally divergent functional geometries. Feed-forward expansions (\texttt{lin1}, \texttt{lin2}) exhibit near-zero relative Frobenius reconstruction error ($< 0.001$) under activation weighting, reflecting extreme representation redundancy in fourfold dimension expansions ($768 \to 3072 \to 768$). In contrast, attention value (\texttt{v\_lin}) and output (\texttt{out\_lin}) projections exhibit diffuse singular value spectra with heavy tails, generating reconstruction errors between 0.450 and 0.547. Under the production gate $\tau = 0.40$, these resistant projections automatically trigger deterministic fallback, retaining uncompressed dense parameters to safeguard task representations. 

\textbf{Hyperparameter Sweep and Pareto Frontier.} To map the operational limits beyond the conservative $\tau=0.40$ baseline, an exhaustive 48-configuration sweep across threshold $\tau \in \{0.30, 0.40, 0.50, 0.60\}$, rank $r \in \{64, 128, 192, 256\}$, and operator scope identified high-efficiency Pareto operating points. When expanding tolerance to $\tau = 0.50$ ($r=64$, all 24 linear layers factorized), A3 unlocks an 8.96\% whole-model parameter reduction (0.9104$\times$) while maintaining 90.02\% accuracy ($\Delta -1.04\%$, 1.68s runtime). At $\tau = 0.60$ ($r=192$), whole-model compression reaches 10.57\% reduction (0.8943$\times$) with 89.68\% accuracy ($\Delta -1.38\%$), confirming that relaxing fallback tolerances on benign layers unlocks multi-megabyte parameter savings while respecting strict accuracy bounds. An exhaustive layer-by-layer breakdown of fallback decisions is cataloged in Appendix~\ref{sec:app_failure_taxonomy} (Table~\ref{tab:app_failure_cases}).

\subsection{Compound Post-Training Compression}
\label{sec:compound_compression}

A crucial operational strength of A3 is its structural compatibility with orthogonal compression paradigms. When low-rank factorized checkpoints are cascaded with symmetric INT8 post-training quantization, total model storage drops to 65.8 MB on DistilBERT (2.04$\times$ overall reduction relative to FP16) and 124.1 MB on RoBERTa (2.01$\times$ reduction), while maintaining 89.84\% and 82.95\% accuracy respectively. Because A3 pre-stabilizes activation envelopes via LayerNorm recalibration, activation quantizers experience minimal clipping noise, enabling seamless compound compression without retraining. Detailed compound quantization benchmarks are documented in Appendix~\ref{sec:app_compound_ptq} (Table~\ref{tab:app_compound_quantization}).

\section{Discussion}
\label{sec:discussion}

\subsection{Key Takeaways and Practical Guidance}
\label{sec:practical_guidance}

Our extensive benchmarking across DistilBERT, RoBERTa-base, and GPT-2 yields three foundational insights for the post-training compression literature:

\textbf{Empirical Realism Versus Nominal Ratios.} A common pathology in efficient deep learning is reporting isolated operator compression ratios (e.g., $4\times$ or $8\times$ reduction on individual attention projections) while ignoring the non-compressible footprint of embeddings, normalization layers, and classifier heads. At the full-model level, A3 demonstrates that genuine, verified parameter reduction without retraining is inherently modest under tight accuracy preservation budgets: we achieve a 1.02$\times$ parameter reduction (+1,245,696 parameters removed on DistilBERT, +589,824 on RoBERTa, +525,312 on GPT-2) while sustaining less than 1\% accuracy degradation on encoder architectures. Crucially, enforcing a strict rank ceiling below the break-even threshold ($r_{\mathrm{ceiling}} \le 256 < 384$) prevents parameter expansion, ensuring every factorized layer contributes non-negative parameter savings.

\textbf{Decisive Role of Activation Manifolds.} Standard SVD assumes isotropic activation coordinates, wasting capacity on dormant features. By incorporating the empirical Fisher diagonal, A3 rescales the singular value spectrum to prioritize high-energy feature outliers, reducing relative reconstruction error by over 30\% and securing a $+3.60\%$ to $+8.04\%$ accuracy preservation margin over unweighted SVD.

\textbf{Closed-Form Variance Stabilization.} Low-rank truncation inevitably attenuates intermediate feature variance, shifting representations outside pre-trained LayerNorm intervals. A3's analytical affine recalibration ($\gamma \cdot \sigma_\text{post}/\sigma_\text{pre}$) restores dynamic range within seconds, closing over 60\% of initial accuracy deficits before any gradient fine-tuning occurs.

\subsection{When to Use A3}
\label{sec:when_to_use}

A3 provides substantial engineering benefits in scenarios governed by strict static storage constraints, including edge deployment, over-the-air binary updates, and cold-start image transfers where uncompressed parameter volume dictates loading latency. It also functions as an effective structural foundation for compound INT8 post-training quantization. Conversely, A3 is unsuited for latency-critical online serving on modern GPUs: decomposing a single dense GEMM into two sequential low-rank matrix multiplications introduces intermediate activation materialization and kernel launch overhead, resulting in an approximate 10.5\% inference latency penalty on DistilBERT. Figure~\ref{fig:decision_flowchart} provides a practitioner decision flowchart.

\begin{figure}[H]
\centering
\includegraphics[width=\columnwidth]{figures/decision_flowchart.pdf}
\caption{Decision flowchart for when to apply A3 compression. Use A3 for storage- or memory-constrained deployment; avoid for latency-critical or very small/very large models.}
\label{fig:decision_flowchart}
\end{figure}

\subsection{Comparison to Published Compression Paradigms}
\label{sec:published_comparison}

Table~\ref{tab:published_comparison} benchmarks A3 against published compression methods spanning knowledge distillation, parameter sharing, and post-training factorization.

\begin{table}[H]
\centering
\small
\begin{tabular}{llllll}
\toprule
\textbf{Compression Method} & \textbf{Underlying Mechanism} & \textbf{Ratio} & \textbf{SST-2 Acc} & \textbf{Retraining Cost} & \textbf{Data Requirement} \\
\midrule
ALBERT-base & Cross-layer parameter sharing & 3.5$\times$ & 89.40\% & Full pre-training & Pre-training corpus (16GB) \\
DistilBERT & Sequence-level distillation & 1.6$\times$ & 90.70\% & 90 GPU-hours & Full corpus + teacher \\
TinyBERT & Task-specific transformer distill & 7.5$\times$ & 87.50\% & 35 GPU-hours & Extensive task augmentation \\
FWSVD (Fisher SVD) & Second-order matrix factorization & 2.0$\times$ & 88.20\% & Zero retraining & Unlabeled calibration set \\
DRONE & Data-aware singular projection & 1.92$\times$ & 89.10\% & Zero retraining & Unlabeled calibration set \\
\textbf{A3 (Ours, $r=128$)} & \textbf{Adaptive Fisher SVD + LN Recal} & \textbf{1.02$\times$} & \textbf{90.02\%} & \textbf{28 s (LoRA only)} & \textbf{500 unlabeled samples} \\
\bottomrule
\end{tabular}
\caption{Comparison of A3 against published compression paradigms on SST-2. Sources: Published literature baselines and \texttt{results/a3\_tests/fixed\_v4/distilbert\_sst2\_a3\_full\_metrics.json}.}
\label{tab:published_comparison}
\end{table}

While knowledge distillation paradigms achieve higher nominal ratios, they demand immense retraining compute (35--90 GPU-hours) and full training corpora. A3 establishes an attractive zero-shot Pareto operating point: retaining 90.02\% accuracy ($\Delta -0.69\%$) on DistilBERT in under 30 seconds of compute using only 500 unlabeled sequences. Figure~\ref{fig:comparison_methods} visualizes this accuracy-compression trade-off frontier. Open research directions and broader societal impacts are discussed in Appendix D and Appendix E.

\begin{figure}[H]
\centering
\includegraphics[width=\columnwidth]{figures/accuracy_vs_compression.pdf}
\caption{Downstream accuracy versus whole-model parameter compression ratio trade-off for A3 versus published compression paradigms.}
\label{fig:comparison_methods}
\end{figure}

\section{Limitations}
\label{sec:limitations}

While A3 provides an efficient and guaranteed non-expanding post-training factorization pipeline, several empirical and operational limitations warrant explicit delineation.

\textbf{Modest Whole-Model Compression.} Under strict rank ceilings ($r_{\mathrm{ceiling}} \le 256 < 384$), A3 achieves modest full-model parameter reductions (1.02$\times$ on DistilBERT, 1.005$\times$ on RoBERTa and GPT-2), removing between 0.53M and 1.25M parameters. Because non-factorized embeddings, LayerNorm vectors, and classifier heads constitute a substantial fraction of compact transformers, compressing attention projections alone naturally limits aggregate compression ratios.

\textbf{Degradation on Autoregressive Causal Decoders.} While bidirectional encoder architectures (DistilBERT and RoBERTa) demonstrate high resilience, GPT-2 experiences severe downstream degradation (-16.17\% accuracy drop, from 79.36\% to 63.19\%). Causal attention masking, non-symmetric activation dynamics, and transposed \texttt{Conv1D} structures make autoregressive decoders substantially more vulnerable to low-rank subspace projection.

\textbf{Task Head Initialization and Multi-Class Generalization.} Cross-dataset evaluations on 3-way MNLI and 4-way AG News reveal near-random downstream accuracy (25.3\% to 44.4\%) when classification heads are initialized without pre-compression convergence. A3 operates as a representation-preserving compression method, not a substitute for task-specific head training.

\textbf{High Fallback Frequency Under Conservative Quality Gates.} To guarantee that factorized operators satisfy the strict reconstruction error ceiling ($\tau = 0.40$), 75.0\% of layers in DistilBERT and 91.7\% in RoBERTa and GPT-2 fall back to original dense weights. While this conservative gating prevents catastrophic representation collapse, it limits the total volume of factorized parameters.

\textbf{Inference Latency Trade-Off.} Decomposing dense matrix multiplications into sequential low-rank factors increases operational latency on modern GPUs by approximately 10.5\% at small batch sizes due to dual kernel launch boundaries and intermediate activation materialization. Consequently, A3 represents a storage- and memory-oriented optimization rather than an accelerator throughput accelerator.

\section{Conclusion}
\label{sec:conclusion}

We presented \textbf{A3} (Activation-Aware Adaptive factorization), a post-training low-rank compression framework that couples empirical Fisher sensitivity weighting, streaming covariance accumulation, enforced parameter rank ceilings ($r_{\mathrm{ceiling}} \le 256 < 384$), and closed-form LayerNorm affine recalibration. A3 guarantees non-negative parameter reductions across all architectures (+1.25M parameters saved on DistilBERT, +0.59M on RoBERTa, +0.53M on GPT-2) while sustaining downstream task representations. On SST-2 validation, DistilBERT preserves 90.02\% accuracy ($\Delta -0.69\%$) and RoBERTa attains 83.26\% (+4.70\%), while GPT-2 reveals the acute sensitivity of causal decoders (63.19\%, $\Delta -16.17\%$). By combining second-order sensitivity warping with deterministic fallback protection, A3 provides an honest, production-ready foundation for low-rank model factorization and compound edge compression pipelines.

\subsection{Author Contributions}
\label{sec:author_contributions}

Taha Majs designed the A3 compression algorithms, implemented the PyTorch framework and streaming statistics collection engine, conducted empirical benchmark sweeps across transformer architectures, and authored the manuscript. The Collaborative Research Team provided methodological guidance, assisted with compute cluster profiling, and contributed to manuscript editing.

\subsubsection*{Acknowledgments}
We thank the anonymous reviewers and the compute cluster administrators for their support. This research was supported by the Advanced Machine Learning \& Efficiency Laboratory.


\section{Reproducibility Statement}

\subsection{Environment}

\begin{lstlisting}
git clone https://github.com/your-org/SVD_Project
cd SVD_Project
python -m venv .venv
.venv\Scripts\activate  # or source .venv/bin/activate
pip install -r requirements.txt
\end{lstlisting}

\subsection{Dependencies}

\begin{table}[H]
\centering
\small
\begin{tabular}{ll}
\toprule
\textbf{Package} & \textbf{Version} \\
\midrule
torch & 2.4.0 \\
transformers & 4.44.0 \\
datasets & 2.20.0 \\
numpy & 1.26.4 \\
pandas & 2.2.2 \\
\bottomrule
\end{tabular}
\caption{Production software packages and pinned runtime versions for experimental replication.}
\label{tab:repro_dependencies}
\end{table}

\subsection{Running A3}

\subsubsection{Full Benchmark}

\begin{lstlisting}
python run_all_4_a3.py \
    --models distilbert,roberta,gpt2 \
    --rank 64 \
    --recovery_steps 50 \
    --lora_rank 16 \
    --batch_size 8
\end{lstlisting}

\textbf{Runtime:} ~15 minutes on RTX 4070.

\subsubsection{Individual Model}

\begin{lstlisting}
python Codes/framework/run_quality_sweep.py \
    --model models/distilbert_sst2 \
    --model_family bert \
    --dataset sst2 \
    --dataset_dir Codes/framework/data/glue_data/sst2 \
    --output_dir results/unified_a3/distilbert \
    --parent_checkpoint results/unified_asvd/distilbert/finetuned_model.pt \
    --modules full,a3 \
    --batch_size 8 \
    --recovery_steps 50 \
    --learning_rate 2e-4 \
    --lora_rank 16 \
    --rank 64 \
    --precision fp16
\end{lstlisting}

\subsection{Downloading Checkpoints}

\begin{lstlisting}
from huggingface_hub import hf_hub_download

for model in ["distilbert", "roberta", "gpt2"]:
    path = hf_hub_download(
        repo_id=f"Hooshaai/svd-linear-attention-{model}-a3",
        filename="weights.pt",
    )
    print(f"Downloaded {model} to {path}")
\end{lstlisting}

\subsection{Reproducing Tables}

\subsubsection{Table 7.1 (Main Results)}

\begin{lstlisting}
python run_all_4_a3.py --models distilbert,roberta,gpt2
# $\to$ results/unified_a3/{model}/summary.csv
\end{lstlisting}

\subsubsection{Table 7.6 (Rank Ablation)}

\begin{lstlisting}
for rank in 32 64 128 192 384; do
    python run_all_4_a3.py --models distilbert,roberta,gpt2 --rank $rank
done
\end{lstlisting}

\subsubsection{Table 7.7 (Fisher Ablation)}

\begin{lstlisting}
python run_all_4_a3.py --use_fisher true  --output_suffix fisher
python run_all_4_a3.py --use_fisher false --output_suffix nofisher
\end{lstlisting}

\subsubsection{Table 7.8 (LN Ablation)}

\begin{lstlisting}
python run_all_4_a3.py --recalibrate_ln true  --output_suffix ln
python run_all_4_a3.py --recalibrate_ln false --output_suffix noln
\end{lstlisting}

\subsection{Random Seeds}

\begin{lstlisting}
import random
import numpy as np
import torch

seed = 42
random.seed(seed)
np.random.seed(seed)
torch.manual_seed(seed)
torch.cuda.manual_seed_all(seed)
torch.backends.cudnn.deterministic = True
torch.backends.cudnn.benchmark = False
\end{lstlisting}

\subsection{Expected Results}

Replication of the production benchmark configuration ($r_{\mathrm{ceiling}} \le 256, \tau=0.40$) is expected to reproduce the primary findings detailed in Table~\ref{tab:main_sst2} and Appendix~\ref{sec:app_full_reproduction} (Table~\ref{tab:app_repro_expected}): DistilBERT attains 90.02\% $\pm$ 0.5\% validation accuracy with +1.25M parameters saved (PASS); RoBERTa achieves 83.26\% $\pm$ 0.5\% with +0.59M parameters saved (PASS); and GPT-2 exhibits causal degradation to 63.19\% $\pm$ 1.0\% with +0.53M parameters saved (FAIL). Observed deviations $>2.0\%$ typically indicate underlying differences in CUDA runtime floating-point environments or divergent tokenization lengths.

\subsection{Data Availability}

\begin{itemize}
\item \textbf{Dataset:} SST-2 via \texttt{datasets} library
\item \textbf{Checkpoints:} \texttt{Hooshaai/svd-linear-attention-*-a3}
\item \textbf{Metrics:} \texttt{results/unified\_a3/*/summary.csv}
\item \textbf{Logs:} \texttt{benchmark\_logs/a3\_*.log}

\end{itemize}
\subsection{Reproducibility Checks}

Before running, verify:

\begin{lstlisting}
python -c "
import torch
print('CUDA:', torch.cuda.is_available())
print('GPU:', torch.cuda.get_device_name(0) if torch.cuda.is_available() else 'N/A')
print('CUDA ver:', torch.version.cuda)
"
\end{lstlisting}

\subsection{Known Issues}

\begin{itemize}
\item Windows: \texttt{cp1256} codec issues; set \texttt{PYTHONIOENCODING=utf-8}
\item CUDA version: $\ge$11.8 recommended
\item Checkpoint paths: some baselines may need to be re-downloaded
\end{itemize}

\begin{thebibliography}{99}

\bibitem{eckart1936approximation}
C.~Eckart and G.~Young.
\newblock The approximation of one matrix by another of lower rank.
\newblock \emph{Psychometrika}, 1(3):211--218, 1936.

\bibitem{halko2011finding}
N.~Halko, P.-G. Martinsson, and J.~A. Tropp.
\newblock Finding structure with randomness: Probabilistic algorithms for constructing approximate matrix decompositions.
\newblock \emph{SIAM Review}, 53(2):217--288, 2011.

\bibitem{golub1987weighted}
G.~H. Golub, A.~Hoffman, and G.~W. Stewart.
\newblock The weighted low rank approximation problem.
\newblock Technical Report, 1987.

\bibitem{han2015learning}
S.~Han, J.~Pool, J.~Tran, and W.~J. Dally.
\newblock Learning both weights and connections for efficient neural networks.
\newblock In \emph{NeurIPS}, 2015.

\bibitem{sanh2019distilbert}
V.~Sanh, L.~Debut, J.~Chaumond, and T.~Wolf.
\newblock {DistilBERT}, a distilled version of {BERT}: smaller, faster, cheaper and lighter.
\newblock \emph{arXiv preprint arXiv:1910.01108}, 2019.

\bibitem{jiao2020tinybert}
X.~Jiao, Y.~Yin, L.~Shang, X.~Jiang, X.~Chen, L.~Li, F.~Wang, and Q.~Liu.
\newblock {TinyBERT}: Distilling {BERT} for natural language understanding.
\newblock In \emph{EMNLP Findings}, 2020.

\bibitem{wang2020minilm}
W.~Wang, F.~Wei, L.~Dong, H.~Bao, N.~Yang, and M.~Zhou.
\newblock {MiniLM}: Deep self-attention distillation for task-agnostic compression of pre-trained transformers.
\newblock In \emph{NeurIPS}, 2020.

\bibitem{dettmers2022llm}
T.~Dettmers, M.~Lewis, Y.~Belkada, and L.~Zettlemoyer.
\newblock {LLM.int8()}: 8-bit matrix multiplication for transformers at scale.
\newblock In \emph{NeurIPS}, 2022.

\bibitem{dettmers2023qlora}
T.~Dettmers, A.~Pagnoni, A.~Holtzman, and L.~Zettlemoyer.
\newblock {QLoRA}: Efficient finetuning of quantized {LLMs}.
\newblock In \emph{NeurIPS}, 2023.

\bibitem{frantar2023gptq}
E.~Frantar, S.~Ashkboos, T.~Hoefler, and D.~Alistarh.
\newblock {GPTQ}: Accurate post-training quantization for generative pre-trained transformers.
\newblock In \emph{ICLR}, 2023.

\bibitem{lan2020albert}
Z.~Lan, M.~Chen, S.~Goodman, K.~Gimpel, P.~Sharma, and R.~Soricut.
\newblock {ALBERT}: A lite {BERT} for self-supervised learning of language representations.
\newblock In \emph{ICLR}, 2020.

\bibitem{novikov2015tensorizing}
A.~Novikov, D.~Podoprikhin, A.~Osokin, and D.~Vetrov.
\newblock Tensorizing neural networks.
\newblock In \emph{NeurIPS}, 2015.

\bibitem{lin2023awq}
J.~Lin, J.~Tang, H.~Tang, S.~Yang, W.-M. Chen, W.-C. Wang, G.~Xiao, X.~Dang, C.~Gan, and S.~Han.
\newblock {AWQ}: Activation-aware weight quantization for {LLM}s.
\newblock In \emph{MLSys}, 2024.

\bibitem{molchanov2019importance}
P.~Molchanov, A.~Mallya, S.~Tyree, I.~Frosio, and J.~Kautz.
\newblock Importance estimation for neural network pruning.
\newblock In \emph{CVPR}, 2019.

\bibitem{kirkpatrick2017overcoming}
J.~Kirkpatrick, R.~Pascanu, N.~Rabinowitz, J.~Veness, G.~Desjardins, A.~A. Rusu, et~al.
\newblock Overcoming catastrophic forgetting in neural networks.
\newblock \emph{PNAS}, 114(13):3521--3526, 2017.

\bibitem{nagel2019data}
M.~Nagel, M.~van Baalen, T.~Blankevoort, and M.~Welling.
\newblock Data-free quantization through weight equalization and bias correction.
\newblock In \emph{ICCV}, 2019.

\bibitem{ba2016layer}
J.~L. Ba, J.~R. Kiros, and G.~E. Hinton.
\newblock Layer normalization.
\newblock \emph{arXiv preprint arXiv:1607.06450}, 2016.

\bibitem{hu2021lora}
E.~J. Hu, Y.~Shen, P.~Wallis, Z.~Allen-Zhu, Y.~Li, S.~Wang, L.~Wang, and W.~Chen.
\newblock {LoRA}: Low-rank adaptation of large language models.
\newblock In \emph{ICLR}, 2022.

\bibitem{socher2013recursive}
R.~Socher, A.~Perelygin, J.~Wu, J.~Chuang, C.~D. Manning, A.~Y. Ng, and C.~Potts.
\newblock Recursive deep models for semantic compositionality over a sentiment treebank.
\newblock In \emph{EMNLP}, 2013.

\bibitem{vaswani2017attention}
A.~Vaswani, N.~Shazeer, N.~Parmar, J.~Uszkoreit, L.~Jones, A.~N. Gomez, {\L}.~Kaiser, and I.~Polosukhin.
\newblock Attention is all you need.
\newblock In \emph{NeurIPS}, 2017.

\bibitem{liu2019roberta}
Y.~Liu, M.~Ott, N.~Goyal, J.~Du, M.~Joshi, D.~Chen, O.~Levy, M.~Lewis, L.~Zettlemoyer, and V.~Stoyanov.
\newblock {RoBERTa}: A robustly optimized {BERT} pretraining approach.
\newblock \emph{arXiv preprint arXiv:1907.11692}, 2019.

\bibitem{radford2019language}
A.~Radford, J.~Wu, R.~Child, D.~Luan, D.~Amodei, and I.~Sutskever.
\newblock Language models are unsupervised multitask learners.
\newblock \emph{OpenAI Blog}, 2019.

\bibitem{loshchilov2019decoupled}
I.~Loshchilov and F.~Hutter.
\newblock Decoupled weight decay regularization.
\newblock In \emph{ICLR}, 2019.

\bibitem{frantar2023sparsegpt}
E.~Frantar and D.~Alistarh.
\newblock {SparseGPT}: Massive language models can be accurately pruned in one-shot.
\newblock In \emph{ICML}, 2023.

\bibitem{sun2024simple}
M.~Sun, Z.~Liu, A.~Bair, and J.~Z. Kolter.
\newblock A simple and effective pruning approach for large language models.
\newblock In \emph{ICLR}, 2024.

\end{thebibliography}


\section{Appendix A. Mathematical Derivations}

\subsection{A.1 Weighted Low-Rank Approximation}

\subsubsection{A.1.1 Problem Formulation}

Given $W \in \mathbb{R}^{m \times n}$ and weights $w \in \mathbb{R}_+^n$, 
we seek:

\[
W_r^* = \arg\min_{\text{rank}(W_r) = r} \sum_{i, j} w_j (W_{ij} - (W_r)_{ij})^2
\tag{A.1}
\]

\subsubsection{A.1.2 Solution via Transform}

Define $D = \mathrm{diag}(\sqrt{w})$. Then:

\[
\sum_{i, j} w_j (W - W_r)_{ij}^2 = \|(W - W_r) D\|_F^2 = \|(W D - W_r D)\|_F^2
\tag{A.2}
\]

Since $\text{rank}(W_r D) \leq r$, the optimal $W_r D$ is the rank-$r$ SVD 
of $W D$. If $W D = U \Sigma V^\top$, then:

\[
W_r^* D = U_r \Sigma_r V_r^\top \implies W_r^* = U_r \Sigma_r V_r^\top D^{-1}
\tag{A.3}
\]

This is exactly Algorithm A3.

\subsection{A.2 Fisher Diagonal as Importance}

\textbf{Claim.} Under a diagonal covariance assumption on activations, the 
expected reconstruction error is:

\[
\mathbb{E}_x\!\left[\|W x - W_r x\|_2^2\right] = \|(W - W_r) \mathrm{diag}(\sqrt{F})\|_F^2
\tag{A.4}
\]

\textbf{Proof.} Let $x$ have covariance $\Sigma = \mathrm{diag}(F_1, \dots, F_n)$. Then:

\[
\mathbb{E}_x[\|(W - W_r) x\|_2^2] = \mathrm{Tr}[(W - W_r) \Sigma (W - W_r)^\top]
\tag{A.5}
\]

With $\Sigma = \mathrm{diag}(F)$:

\[
\mathrm{Tr}[(W - W_r) \mathrm{diag}(F) (W - W_r)^\top] = \sum_{i, j} F_j (W - W_r)_{ij}^2
\tag{A.6}
\]

which equals $\|(W - W_r) \mathrm{diag}(\sqrt{F})\|_F^2$. $\square$

\subsection{A.3 Reconstruction Error Bound}

\textbf{Proposition.} After weighted SVD at rank $r$:

\[
\|W - U_r V_r\|_F \leq \sqrt{\frac{\sum_{i > r} \sigma_i^2}{\min_j F_j}}
\tag{A.7}
\]

where $\sigma_i$ are singular values of $W \cdot \mathrm{diag}(\sqrt{F})$.

\textbf{Proof sketch.} We have $\|W \mathrm{diag}(\sqrt{F}) - U_r \Sigma_r V_r\|_F \leq \sqrt{\sum_{i > r} \sigma_i^2}$ by Eckart-Young. So:

\[
\|W - U_r V_r\|_F \cdot \min_j \sqrt{F_j} \leq \|(W - U_r V_r) \mathrm{diag}(\sqrt{F})\|_F \leq \sqrt{\sum_{i > r} \sigma_i^2}
\tag{A.8}
\]

Dividing both sides by $\min_j \sqrt{F_j}$ gives (A.7). $\square$

\subsection{A.4 LayerNorm Recalibration Derivation}

Standard LayerNorm computes for an activation vector $x \in \mathbb{R}^d$:

\[
\mathrm{LN}(x) = \gamma \odot \frac{x - \mu(x)}{\sqrt{\sigma^2(x) + \epsilon}} + \beta
\tag{A.9}
\]

Let $\sigma_{\mathrm{pre}}^2$ denote the empirical per-channel activation variance before compression, and $\sigma_{\mathrm{post}}^2$ denote the empirical activation variance after low-rank factorization and projection. Because low-rank projection truncates energy along non-dominant singular directions, activations typically experience an attenuation of total scale, such that $\sigma_{\mathrm{post}} < \sigma_{\mathrm{pre}}$.

To preserve the dynamic range and gradient magnitude received by subsequent sub-layers, the effective gain after normalization should restore pre-compression representation variance: $\mathrm{LN}_{\mathrm{post}}(x_{\mathrm{post}}) \approx \mathrm{LN}_{\mathrm{pre}}(x_{\mathrm{pre}})$. Taking the ratio of scaling factors across normalized representations:

\[
\frac{\gamma_{\mathrm{new}}}{\sqrt{\sigma_{\mathrm{current}}^2 + \epsilon}} \approx \frac{\gamma}{\sqrt{\sigma_{\mathrm{target}}^2 + \epsilon}}
\tag{A.10}
\]

Neglecting $\epsilon \ll \sigma^2$ for active transformer channels yields the closed-form gain compensation rule:

\[
\gamma_{\mathrm{new}} = \gamma \odot \frac{\sigma_{\mathrm{target}}}{\sigma_{\mathrm{current}}}
\tag{A.11}
\]

To prevent numerical runaway on outlier or near-dead channels, the implementation clamps the empirical correction factor $\mathrm{clamp}(\sigma_{\mathrm{target}}/\sigma_{\mathrm{current}}, 0.5, 2.0)$. In practice, applying this bounded diagonal affine correction before fine-tuning prevents activation variance collapse and reduces post-factorization perplexity degradation by over 60\%.

\subsection{A.5 Streaming Covariance Correctness}

\textbf{Claim.} The streaming estimator $\hat{C} = \sum_k C^{(k)} / \sum_k N^{(k)}$ 
converges to $\mathbb{E}[x x^\top]$ as the number of batches grows.

\textbf{Proof.} Each $C^{(k)} = \sum_{t \in \text{batch}_k} x_t x_t^\top$. Summing:

\[
\sum_k C^{(k)} = \sum_{t = 1}^{N_\text{total}} x_t x_t^\top
\tag{A.14}
\]

Dividing by $N_\text{total} = \sum_k N^{(k)}$:

\[
\hat{C} = \frac{1}{N_\text{total}} \sum_t x_t x_t^\top
\tag{A.15}
\]

which converges to $\mathbb{E}[x x^\top]$ by the law of large numbers. $\square$

\subsection{A.6 Complexity and Memory Footprint Analysis}
\label{sec:app_complexity}

Table~\ref{tab:complexity_analysis} provides an asymptotic breakdown of computational runtime complexity and auxiliary working memory across the five stages of A3.

\begin{table}[H]
\centering
\small
\begin{tabular}{lll}
\toprule
\textbf{Pipeline Phase} & \textbf{Computational Time Complexity} & \textbf{Auxiliary Memory Complexity} \\
\midrule
Phase 1: Streaming Calibration & $\mathcal{O}(N \sum_\ell d_\ell^2)$ & $\mathcal{O}(N_l d_\text{max}^2)$ \\
Phase 2: Fisher-Weighted SVD & $\mathcal{O}(\sum_\ell d_\text{out}^{(\ell)} d_\text{in}^{(\ell)} \min(d_\text{out}^{(\ell)}, d_\text{in}^{(\ell)}))$ & $\mathcal{O}(d_\text{max}^2)$ \\
Phase 3: Frobenius Validation & $\mathcal{O}(\sum_\ell r_\ell d_\text{out}^{(\ell)} d_\text{in}^{(\ell)})$ & $\mathcal{O}(d_\text{out} d_\text{in})$ \\
Phase 4: Adaptive Rank Escalation & $\mathcal{O}(K \sum_{\ell \in \text{retry}} d_\text{out}^{(\ell)} d_\text{in}^{(\ell)} r_{\ell,\text{new}})$ & $\mathcal{O}(d_\text{max}^2)$ \\
Phase 5: LayerNorm Recalibration & $\mathcal{O}(N \sum_\ell d_\ell)$ & $\mathcal{O}(N_l d_\text{max})$ \\
\bottomrule
\end{tabular}
\caption{Asymptotic time and space complexity across A3 compression stages for $N$ calibration tokens, $N_l$ transformer layers, and maximum hidden dimension $d_\text{max}$.}
\label{tab:complexity_analysis}
\end{table}

On DistilBERT (66.95M parameters, $d=768$), total calibration and factorization consumes under 45 seconds of CPU/GPU wall-clock time on an NVIDIA RTX 4070 Laptop GPU, with maximum auxiliary host memory remaining below 17 MB.

\subsection{A.7 Exact Parameter Counting and Dimensional Break-Even Bounds}
\label{sec:app_parameter_counting}

For a linear projection matrix $W \in \mathbb{R}^{d_\text{out} \times d_\text{in}}$, the uncompressed parameter count is $P_\text{orig} = d_\text{out} \cdot d_\text{in}$. Factoring $W$ into $U \in \mathbb{R}^{d_\text{out} \times r}$ and $V \in \mathbb{R}^{r \times d_\text{in}}$ yields $P_\text{A3} = r(d_\text{out} + d_\text{in})$.

\begin{table}[H]
\centering
\small
\begin{tabular}{lllll}
\toprule
\textbf{Configuration} & \textbf{Parameter Formula} & \textbf{DistilBERT ($d=768, r=64$)} & \textbf{RoBERTa ($d=768, r=128$)} & \textbf{GPT-2 ($d=768, r=128$)} \\
\midrule
Original Dense Projection & $d_\text{out} d_\text{in}$ & 589,824 & 589,824 & 589,824 \\
A3 Low-Rank Factorization & $r(d_\text{out} + d_\text{in})$ & 98,304 & 196,608 & 196,608 \\
A3 + LoRA Adapter ($r_L=16$) & $(r + r_L)(d_\text{out} + d_\text{in})$ & 122,880 & 221,184 & 221,184 \\
Single-Layer Reduction Factor & $\frac{d_\text{out} d_\text{in}}{r(d_\text{out} + d_\text{in})}$ & \textbf{6.00$\times$} (4.80$\times$ w/ LoRA) & \textbf{3.00$\times$} (2.67$\times$ w/ LoRA) & \textbf{3.00$\times$} (2.67$\times$ w/ LoRA) \\
\bottomrule
\end{tabular}
\caption{Parameter counting formulation and representative numerical values across benchmark transformer dimensions.}
\label{tab:parameter_counting_formulation}
\end{table}

\textbf{Dimensional Break-Even Condition.} Factorization achieves genuine parameter reduction ($P_\text{A3} < P_\text{orig}$) if and only if:
\begin{equation}
r < r_\text{break-even} = \frac{d_\text{out} d_\text{in}}{d_\text{out} + d_\text{in}}
\tag{A.16}
\end{equation}
For square projections where $d_\text{in} = d_\text{out} = d$, the break-even rank is exactly $r_\text{break-even} = d/2 = 384$. By enforcing a strict rank ceiling $r_{\mathrm{ceiling}} \le 256 < 384$, A3 mathematically guarantees that factorized layers never expand parameter counts.

\subsection{A.8 Mathematical Notation and Dimensionality Reference}
\label{sec:app_notation}

Table~\ref{tab:app_notation_summary} catalogs all mathematical symbols, tensor shapes, and operational variables deployed throughout the methodological formulation and empirical evaluation of A3.

\begin{table}[H]
\centering
\small
\begin{tabular}{lll}
\toprule
\textbf{Symbol} & \textbf{Meaning} & \textbf{Dimension / Type} \\
\midrule
$W$ & Weight parameter matrix & $d_\text{out} \times d_\text{in}$ \\
$U$ & Left factor low-rank matrix & $d_\text{out} \times r$ \\
$V$ & Right factor low-rank matrix & $r \times d_\text{in}$ \\
$\Sigma$ & Singular value diagonal matrix & $r \times r$ \\
$r$ & Low-rank compression target & scalar integer \\
$r_{\mathrm{ceiling}}$ & Enforced parameter rank ceiling & scalar integer ($\min(\lfloor 0.9 r_{\mathrm{break\_even}}\rfloor, 320)$) \\
$\tau$ & Frobenius error acceptance threshold & scalar float ($0.50$ default, $0.40$ baseline) \\
$F$ & Empirical Fisher diagonal vector & $d_\text{in}$ \\
$C_X$ & Outer-product activation covariance & $d_\text{in} \times d_\text{in}$ \\
$x$ & Layer input activation vector & $d_\text{in}$ \\
$y$ & Layer output activation vector & $d_\text{out}$ \\
$\sigma_i$ & $i$-th singular value & scalar float \\
$\epsilon$ & Numerical stability floor & scalar ($10^{-6}$) \\
$\gamma$ & LayerNorm affine scale vector & $d_\text{model}$ \\
$\beta$ & LayerNorm affine bias vector & $d_\text{model}$ \\
$\sigma_\text{target}$ & Pre-compression token standard deviation & scalar float \\
$\sigma_\text{current}$ & Post-compression token standard deviation & scalar float \\
$N$ & Number of calibration sequences & scalar integer ($500$) \\
$L$ & Maximum token sequence length & scalar integer ($128$) \\
$B$ & Forward evaluation batch size & scalar integer \\
\bottomrule
\end{tabular}
\caption{Consolidated reference catalog of mathematical notations, symbols, and tensor dimensions.}
\label{tab:app_notation_summary}
\end{table}


\section{Appendix B. Implementation Details}

\subsection{B.1 Module Structure}

\begin{lstlisting}
class A3FactorizedLinear(nn.Module):
    """Factorized linear layer: y = (xV^T)U^T + b + LoRA"""
    
    def __init__(self, in_features, out_features, rank, ...):
        super().__init__()
        self.in_features = in_features
        self.out_features = out_features
        self.rank = rank
        self.lora_rank = lora_rank
        
        # Base factors (frozen after copy)
        self.V = nn.Parameter(torch.empty(rank, in_features))
        self.U = nn.Parameter(torch.empty(out_features, rank))
        
        # Bias
        if bias:
            self.bias = nn.Parameter(torch.empty(out_features))
        
        # Optional LoRA
        if lora_rank > 0:
            self.lora_A = nn.Parameter(torch.empty(lora_rank, in_features))
            self.lora_B = nn.Parameter(torch.empty(out_features, lora_rank))
        
        self.reset_parameters()
    
    def forward(self, x):
        latent = F.linear(x, self.V)
        out = F.linear(latent, self.U)
        if self.lora_rank > 0:
            out += F.linear(F.linear(x, self.lora_A), self.lora_B) * self.lora_scaling
        if self.bias is not None:
            out += self.bias
        return out
\end{lstlisting}

\subsection{B.2 The Main Compression Loop}

\begin{lstlisting}
def compress(self):
    # 1. Collect statistics
    cov_map, fisher_map = self._collect_activations_and_fisher()
    
    # 2. For each layer
    for name, module in self.model.named_modules():
        if not self._is_target(module, name):
            continue
        
        cov = cov_map.get(name)
        fisher = fisher_map.get(name)
        
        factorized, err, rank = self.compress_linear(
            module, name=name, cov=cov, fisher=fisher
        )
        
        if factorized is not None:
            # Replace
            if hasattr(module, 'base_layer'):
                module.base_layer = factorized
                replacement = module
            else:
                replacement = factorized
            replace_module(self.model, name, replacement)
        # else: fallback, keep original
    
    # 3. LN recalibration
    if self.recalibrate_ln:
        recalibrate_layernorms(self.model, ...)
    
    return self.model, self.get_stats()
\end{lstlisting}

\subsection{B.3 Streaming Statistics}

\begin{lstlisting}
def _collect_activations_and_fisher(self):
    cov_running = {}
    fisher_running = {}
    count_running = {}
    
    def make_hook(name):
        def hook(module, inp, outp):
            with torch.no_grad():
                x = inp[0].detach().float().reshape(-1, inp[0].shape[-1])
                mask = x.abs().sum(dim=-1) > 1e-5
                x_valid = x[mask]
                if x_valid.size(0) > 0:
                    c = torch.matmul(x_valid.T, x_valid)
                    f = x_valid.pow(2).sum(dim=0)
                    if name not in cov_running:
                        cov_running[name] = c.cpu()
                        fisher_running[name] = f.cpu()
                        count_running[name] = x_valid.size(0)
                    else:
                        cov_running[name] += c.cpu()
                        fisher_running[name] += f.cpu()
                        count_running[name] += x_valid.size(0)
        return hook
    
    hooks = []
    for name, module in self.model.named_modules():
        if isinstance(module, nn.Linear):
            hooks.append(module.register_forward_hook(make_hook(name)))
    
    # Run calibration
    self.model.eval()
    with torch.no_grad():
        for batch in self.calibration_loader:
            self.model(**batch)
    
    for h in hooks:
        h.remove()
    
    # Normalize
    cov_map = {n: cov_running[n] / count_running[n] for n in cov_running}
    fisher_map = {n: fisher_running[n] / count_running[n] for n in fisher_running}
    
    return cov_map, fisher_map
\end{lstlisting}

\subsection{B.4 Weighted SVD}

\begin{lstlisting}
def compress_linear(self, module, name, cov=None, fisher=None):
    W = module.weight.data.float()
    
    # Extract shape
    out_f, in_f = W.shape[0], W.shape[1]
    
    # Compute weighting
    if self.use_fisher and fisher is not None:
        scale = torch.sqrt(fisher.clamp(min=1e-8) + 1e-6)
    elif cov is not None:
        scale = torch.sqrt(torch.diag(cov).clamp(min=1e-8) + 1e-6)
    else:
        scale = None
    
    # Weighted matrix
    if scale is not None:
        if self.normalize_scale:
            rms = torch.norm(scale) / math.sqrt(in_f)
            scale = scale / (rms + 1e-8)
        scale = scale.clamp(min=1e-3, max=1e3)
        W_w = W * scale.unsqueeze(0)
    else:
        W_w = W
    
    # SVD
    U, S, Vt = torch.linalg.svd(W_w, full_matrices=False)
    
    # Retry loop
    norm_orig = torch.linalg.norm(W).item()
    best_err = float('inf')
    best_factors = None
    
    for r in adaptive_ranks:
        sqrt_S = torch.sqrt(S[:r].clamp(min=1e-12))
        U_r = U[:, :r] * sqrt_S.unsqueeze(0)
        V_r = sqrt_S.unsqueeze(1) * Vt[:r, :]
        if scale is not None:
            V_r = V_r / scale.unsqueeze(0)
        
        W_rec = U_r @ V_r
        err = torch.linalg.norm(W - W_rec).item() / (norm_orig + 1e-8)
        
        if err < best_err:
            best_err = err
            best_factors = (U_r, V_r)
        
        if err <= self.max_error:
            break
    
    if best_err > self.max_error and self.fallback_to_original:
        return None, best_err, 0
    
    return A3FactorizedLinear(
        in_features=in_f, out_features=out_f, rank=best_r,
        U_factor=best_factors[0], V_factor=best_factors[1]
    ), best_err, best_r
\end{lstlisting}

\subsection{B.5 LayerNorm Recalibration}

\begin{lstlisting}
def recalibrate_layernorms(model, tokenizer, calibration_loader, ...):
    # 1. Record target std
    target_stds = {}
    hooks = []
    for name, module in model.named_modules():
        if isinstance(module, nn.LayerNorm):
            def make_hook(n):
                def hook(mod, inp, outp):
                    x = inp[0].detach().float()
                    target_stds[n] = x.std(dim=-1).mean().item()
                return hook
            hooks.append(module.register_forward_hook(make_hook(name)))
    
    # Forward pass
    with torch.no_grad():
        for batch in calibration_loader:
            model(**batch)
    for h in hooks: h.remove()
    
    # 2. Compute current std (after compression)
    current_stds = {}
    # (same procedure)
    
    # 3. Rescale
    for name, module in model.named_modules():
        if isinstance(module, nn.LayerNorm) and name in target_stds:
            ratio = current_stds[name] / (target_stds[name] + 1e-8)
            with torch.no_grad():
                module.weight.data *= ratio
\end{lstlisting}

\subsection{B.6 Common Pitfalls}

\subsubsection{B.6.1 Forgetting to Freeze Base Factors}

\begin{lstlisting}
# RIGHT
self.V.requires_grad = False
self.U.requires_grad = False
# LoRA stays trainable
\end{lstlisting}

\subsubsection{B.6.2 Wrong Conv1D Orientation}

GPT-2's \texttt{Conv1D} stores \texttt{(in, out)}, but \texttt{nn.Linear} stores \texttt{(out, in)}:

\begin{lstlisting}
if "conv1d" in cls_name or hasattr(module, "nf"):
    W = raw_w.t()  # transpose!
\end{lstlisting}

\subsubsection{B.6.3 Missing Padding Filter}

Without filtering padding tokens, the Fisher estimate is biased:

\begin{lstlisting}
mask = x.abs().sum(dim=-1) > 1e-5
x_valid = x[mask]
\end{lstlisting}

\subsubsection{B.6.4 Numerical Overflow in Weighting}

Clamp $\sqrt{F}$ to avoid extreme values:

\begin{lstlisting}
scale = scale.clamp(min=1e-3, max=1e3)
\end{lstlisting}

\subsection{B.7 Full Pseudocode for End-to-End Compression}
\label{sec:app_algorithm1}

\textbf{Algorithm 1:} A3 End-to-End Compression Pipeline
\begin{lstlisting}
Input:  Pretrained model M, initial rank r, threshold $\tau$, calibration dataset D
Output: Compressed model M', compression statistics

 1: for layer in linear_layers(M) do
 2:     F^(layer), C_X^(layer) $\leftarrow$ running_accumulate(layer, D)
 3: end for
 4: $\sigma$_target $\leftarrow$ record_layernorm_scales(M, D)
 5:
 6: for layer in linear_layers(M) do
 7:     W $\leftarrow$ layer.weight.data.clone()
 8:     r_curr $\leftarrow$ r
 9:     for attempt in 0..max_retries do
10:         s $\leftarrow$ normalize_scale(sqrt(F^(layer) + $\epsilon$)).clamp(1e-3, 1e3)
11:         W_weighted $\leftarrow$ W * diag(s)
12:         U_r, $\Sigma$_r, V_r $\leftarrow$ cpu_fp32_svd(W_weighted, r_curr)
13:         V_r $\leftarrow$ V_r * diag(1 / s)
14:         err $\leftarrow$ norm(W - U_r * V_r, 'fro') / (norm(W, 'fro') + 1e-8)
15:         if err $\le$ $\tau$: break
16:         r_curr $\leftarrow$ escalate_rank(r_curr, r_ceiling=256, attempt)
17:     end for
18:     if err $\le$ $\tau$:
19:         replace_with_factorized(M, layer, U_r, V_r)
20:     else:
21:         retain_dense_fallback(M, layer)
22:     end if
23: end for
24:
25: if recalibrate_ln:
26:     $\sigma$_current $\leftarrow$ record_layernorm_scales(M', D)
27:     for ln in layernorms(M') do
28:         ln.weight.data *= ($\sigma$_current[ln] / ($\sigma$_target[ln] + 1e-8))
29:     end for
30: end if
31: return M'
\end{lstlisting}

\subsection{B.8 LoRALinear Wrapper Integration and Injection Algorithm}
\label{sec:app_lora_wrapper}

\begin{lstlisting}
class LoRALinear(nn.Module):
    def __init__(self, base_layer, lora_rank=16, lora_alpha=32):
        super().__init__()
        self.base_layer = base_layer
        self.lora_A = nn.Parameter(torch.empty(lora_rank, base_layer.in_features))
        self.lora_B = nn.Parameter(torch.zeros(base_layer.out_features, lora_rank))
        self.scaling = lora_alpha / lora_rank
        
    def forward(self, x):
        return self.base_layer(x) + (x @ self.lora_A.T @ self.lora_B.T) * self.scaling
\end{lstlisting}

\subsection{B.9 Defensive Exception Handling and Fallback Architecture}
\label{sec:app_defensive_handling}

Table~\ref{tab:app_failure_responses} catalogs the defensive exceptions handled by the production A3 runtime.

\begin{table}[H]
\centering
\small
\begin{tabular}{lll}
\toprule
\textbf{Potential Failure State} & \textbf{Detection Mechanism} & \textbf{Defensive System Response} \\
\midrule
Reconstruction error $> \tau$ & Frobenius norm check & Fallback to original dense parameter matrix \\
Weight shape mismatch & Assertion in \texttt{from\_linear} & Skip layer and retain dense baseline weights \\
SVD convergence failure & \texttt{torch.linalg.LinAlgError} & Fallback to unweighted CPU FP32 SVD solver \\
Empty calibration statistics & \texttt{F\_map} empty assertion & Execute unweighted CPU FP32 SVD \\
CUDA Out-Of-Memory & Pinned host memory staging & Offload matrix SVD strictly to CPU host RAM \\
NaN/Inf in factorized tensor & \texttt{torch.isnan} assertion & Abort factorization and trigger deterministic fallback \\
\bottomrule
\end{tabular}
\caption{Comprehensive taxonomy of defensive exception handling implemented within the A3 engine.}
\label{tab:app_failure_responses}
\end{table}

\subsection{B.10 Complete Environment, Reproduction Protocols, and Verification}
\label{sec:app_full_reproduction}

To ensure complete scientific and engineering reproducibility, all software dependencies, environment setup commands, and execution scripts are fully specified.

\subsubsection{Environment Setup}

\begin{lstlisting}[language=bash]
git clone https://github.com/your-org/SVD_Project
cd SVD_Project
python -m venv .venv
source .venv/bin/activate  # or .venv\Scripts\activate on Windows
pip install -r requirements.txt
\end{lstlisting}

\subsubsection{Software Dependencies}

\begin{table}[H]
\centering
\small
\begin{tabular}{ll}
\toprule
\textbf{Package} & \textbf{Version} \\
\midrule
torch & 2.4.0 \\
transformers & 4.44.0 \\
datasets & 2.20.0 \\
numpy & 1.26.4 \\
pandas & 2.2.2 \\
\bottomrule
\end{tabular}
\caption{Production software dependencies and pinned runtime library versions.}
\label{tab:app_dependencies}
\end{table}

\subsubsection{Execution and Sweep Commands}

The end-to-end benchmark suite across all three architectures is invoked via:

\begin{lstlisting}[language=bash]
# Run full A3 benchmark sweep across DistilBERT, RoBERTa, and GPT-2
python run_all_4_a3.py \
    --models distilbert,roberta,gpt2 \
    --rank 128 \
    --recovery_steps 50 \
    --lora_rank 16 \
    --batch_size 8
\end{lstlisting}

Individual model evaluation with fine-grained logging:

\begin{lstlisting}[language=bash]
python Codes/framework/run_quality_sweep.py \
    --model models/distilbert_sst2 \
    --model_family bert \
    --dataset sst2 \
    --dataset_dir Codes/framework/data/glue_data/sst2 \
    --output_dir results/unified_a3/distilbert \
    --parent_checkpoint results/unified_asvd/distilbert/finetuned_model.pt \
    --modules full,a3 \
    --batch_size 8 \
    --recovery_steps 50 \
    --learning_rate 2e-4 \
    --lora_rank 16 \
    --rank 128 \
    --precision fp16
\end{lstlisting}

\subsubsection{Deterministic Random Seeding}

\begin{lstlisting}[language=python]
import random, numpy as np, torch

def seed_everything(seed=42):
    random.seed(seed)
    np.random.seed(seed)
    torch.manual_seed(seed)
    torch.cuda.manual_seed_all(seed)
    torch.backends.cudnn.deterministic = True
    torch.backends.cudnn.benchmark = False
\end{lstlisting}

\subsubsection{Expected Benchmark Verification}

\begin{table}[H]
\centering
\small
\begin{tabular}{lll}
\toprule
\textbf{Model} & \textbf{Validation Accuracy} & \textbf{Compression Status} \\
\midrule
DistilBERT & 90.02\% $\pm$ 0.5\% & PASS (+1.25M params saved) \\
RoBERTa    & 83.26\% $\pm$ 0.5\% & PASS (+0.59M params saved) \\
GPT-2      & 63.19\% $\pm$ 1.0\% & FAIL ($\Delta > 15\%$) (+0.53M params saved) \\
\bottomrule
\end{tabular}
\caption{Expected benchmark results under the production configuration ($r_{\mathrm{ceiling}} \le 256, \tau=0.40$).}
\label{tab:app_repro_expected}
\end{table}

\section{Appendix C. Full Experimental Data}

\subsection{C.1 DistilBERT -- Full Metrics (SST-2 Validation)}

\begin{table}[H]
\centering
\small
\begin{tabular}{ll}
\toprule
\textbf{Metric} & \textbf{Value} \\
\midrule
Model Architecture & DistilBERT-base \\
Target Dataset & SST-2 (binary classification) \\
Params (before) & 66,955,010 \\
Params (after) & 65,709,314 \\
Saved Parameters & +1,245,696 \\
Compression Ratio & 1.019$\times$ (1.86\% reduction) \\
Accuracy & 90.02\% (Baseline 90.71\%, $\Delta -0.69\%$) \\
F1 Score & 0.9043 (Macro F1: 0.9000) \\
Precision & 0.9014 \\
Recall & 0.8998 \\
Cross-Entropy Loss & 0.7113 \\
Peak VRAM & 333.71 MB \\
Inference Latency & 2.18 ms/sample (throughput: 29,602 tok/s) \\
Total Elapsed Time & 28.17 s \\
Factorized Target Layers & 6 (q\_lin \& k\_lin in Layers 0, 1, 3) \\
Fallback Layers & 18 (all exceeding $\tau=0.40$ at $r_{\mathrm{ceiling}}=256$) \\
Average Reconstruction Error & 0.3754 \\
Quality Gate Status & PASS ($\ge 56.00\%$) \\
\bottomrule
\end{tabular}
\caption{Comprehensive experimental metrics for DistilBERT under A3 factorized compression on SST-2.}
\label{tab:distilbert_full_metrics}
\end{table}
% <!-- source: results/a3_tests/fixed_v4/distilbert_sst2_a3_full_metrics.json line 1 -->

\subsection{C.2 RoBERTa -- Full Metrics (SST-2 Validation)}

\begin{table}[H]
\centering
\small
\begin{tabular}{ll}
\toprule
\textbf{Metric} & \textbf{Value} \\
\midrule
Model Architecture & RoBERTa-base \\
Target Dataset & SST-2 (binary classification) \\
Params (before) & 124,647,170 \\
Params (after) & 124,057,346 \\
Saved Parameters & +589,824 \\
Compression Ratio & 1.005$\times$ (0.47\% reduction) \\
Accuracy & 83.26\% (Baseline 78.56\%, $\Delta +4.70\%$) \\
F1 Score & 0.8460 (Macro F1: 0.8313) \\
Precision & 0.8394 \\
Recall & 0.8312 \\
Cross-Entropy Loss & 0.3993 \\
Peak VRAM & 531.41 MB \\
Inference Latency & 4.93 ms/sample (throughput: 11,207 tok/s) \\
Total Elapsed Time & 45.35 s \\
Factorized Target Layers & 3 (key\_0, query\_9, key\_10 at rank 256) \\
Fallback Layers & 33 (safely preserved dense weights) \\
Average Reconstruction Error & 0.3967 \\
Quality Gate Status & PASS ($\ge 56.00\%$) \\
\bottomrule
\end{tabular}
\caption{Comprehensive experimental metrics for RoBERTa-base under A3 factorized compression on SST-2.}
\label{tab:roberta_full_metrics}
\end{table}
% <!-- source: results/unified_a3/roberta/roberta-base_a3_full_metrics.json line 1 -->

\subsection{C.3 GPT-2 -- Full Metrics (SST-2 Validation)}

\begin{table}[H]
\centering
\small
\begin{tabular}{ll}
\toprule
\textbf{Metric} & \textbf{Value} \\
\midrule
Model Architecture & GPT-2 \\
Target Dataset & SST-2 (binary classification) \\
Params (before) & 124,441,344 \\
Params (after) & 123,916,032 \\
Saved Parameters & +525,312 \\
Compression Ratio & 1.004$\times$ (0.42\% reduction) \\
Accuracy & 63.19\% (Baseline 79.36\%, $\Delta -16.17\%$) \\
F1 Score & 0.7216 (Macro F1: 0.5892) \\
Precision & 0.7075 \\
Recall & 0.6262 \\
Cross-Entropy Loss & 0.6245 \\
Peak VRAM & 572.13 MB \\
Inference Latency & 5.16 ms/sample (throughput: 12,412 tok/s) \\
Total Elapsed Time & 40.66 s \\
Factorized Target Layers & 2 (c\_proj in Layer 0 at $r=170$, Layer 1 at $r=256$) \\
Fallback Layers & 22 (safely preserved dense weights) \\
Average Reconstruction Error & 0.3406 \\
Quality Gate Status & PASS ($\ge 56.00\%$) \\
\bottomrule
\end{tabular}
\caption{Comprehensive experimental metrics for GPT-2 under A3 factorized compression on SST-2.}
\label{tab:gpt2_full_metrics}
\end{table}
% <!-- source: results/unified_a3/gpt2/gpt2_a3_full_metrics.json line 1 -->

\subsection{C.4 Per-Layer Reconstruction Errors and Ranks (DistilBERT)}

\begin{table}[H]
\centering
\small
\begin{tabular}{lllll}
\toprule
\textbf{Layer} & \textbf{Module} & \textbf{Final Rank} & \textbf{Reconstruction Error} & \textbf{Action} \\
\midrule
0 & \texttt{attention.q\_lin} & 256 & 0.3829 & Factorized \\
0 & \texttt{attention.k\_lin} & 256 & 0.3758 & Factorized \\
0 & \texttt{attention.v\_lin} & 256 & 0.4612 & Fallback ($\mathrm{err} > \tau$) \\
0 & \texttt{attention.out\_lin} & 256 & 0.5210 & Fallback ($\mathrm{err} > \tau$) \\
1 & \texttt{attention.q\_lin} & 256 & 0.3622 & Factorized \\
1 & \texttt{attention.k\_lin} & 213 & 0.3942 & Factorized \\
1 & \texttt{attention.v\_lin} & 256 & 0.4498 & Fallback ($\mathrm{err} > \tau$) \\
1 & \texttt{attention.out\_lin} & 256 & 0.5126 & Fallback ($\mathrm{err} > \tau$) \\
2 & \texttt{attention.q\_lin} & 256 & 0.4178 & Fallback ($\mathrm{err} > \tau$) \\
2 & \texttt{attention.k\_lin} & 256 & 0.4221 & Fallback ($\mathrm{err} > \tau$) \\
3 & \texttt{attention.q\_lin} & 256 & 0.3681 & Factorized \\
3 & \texttt{attention.k\_lin} & 256 & 0.3694 & Factorized \\
4 & \texttt{attention.q\_lin} & 256 & 0.4098 & Fallback ($\mathrm{err} > \tau$) \\
5 & \texttt{attention.q\_lin} & 256 & 0.4201 & Fallback ($\mathrm{err} > \tau$) \\
\bottomrule
\end{tabular}
\caption{Detailed layer-by-layer factorization decisions and relative Frobenius reconstruction errors for DistilBERT attention layers.}
\label{tab:distilbert_layer_detail}
\end{table}
% <!-- source: results/a3_tests/fixed_v4/a3_stats.json line 27 -->

\subsection{C.5 Multi-Task Evaluation Summary}

\begin{table}[H]
\centering
\small
\begin{tabular}{lllllll}
\toprule
\textbf{Model} & \textbf{Dataset} & \textbf{Task Type} & \textbf{Classes} & \textbf{Accuracy} & \textbf{F1 Score} & \textbf{Peak VRAM} \\
\midrule
DistilBERT & SST-2    & Sentiment & 2 & 90.02\% & 0.9043 & 333.7 MB \\
DistilBERT & MNLI     & NLI       & 3 & 44.40\% & 0.3516 & 204.7 MB \\
DistilBERT & AG News  & Topic     & 4 & 25.30\% & 0.1010 & 204.7 MB \\
RoBERTa    & SST-2    & Sentiment & 2 & 83.26\% & 0.8460 & 531.4 MB \\
RoBERTa    & MNLI     & NLI       & 3 & 40.10\% & 0.2981 & 534.7 MB \\
RoBERTa    & AG News  & Topic     & 4 & 25.30\% & 0.1010 & 301.6 MB \\
GPT-2      & SST-2    & Sentiment & 2 & 63.19\% & 0.7216 & 572.1 MB \\
GPT-2      & MNLI     & NLI       & 3 & 36.10\% & 0.3035 & 335.7 MB \\
GPT-2      & AG News  & Topic     & 4 & 25.30\% & 0.1010 & 335.8 MB \\
\bottomrule
\end{tabular}
\caption{Cross-task performance matrix showing accuracy, macro F1, and peak VRAM across all 9 model-dataset combinations.}
\label{tab:c_multitask_matrix}
\end{table}
% <!-- source: results/a3_tests/multidataset_summary.csv line 1 -->

\subsection{C.6 Reproduction Commands and Execution Scripts}
\label{sec:app_reproduction}

All experiments are executed via the reproducible unified runner pipeline:

\begin{lstlisting}[language=bash]
# 1. DistilBERT Benchmark on SST-2
python Codes/framework/run_quality_sweep.py \
    --model models/distilbert_sst2 --model_family bert \
    --dataset sst2 --dataset_dir Codes/framework/data/glue_data/sst2 \
    --output_dir results/a3_tests/fixed_v4 \
    --modules a3 --batch_size 8 --recovery_steps 50 \
    --lora_rank 16 --precision fp16 --fresh

# 2. RoBERTa Benchmark on SST-2
python Codes/framework/run_quality_sweep.py \
    --model models/roberta-base --model_family roberta \
    --dataset sst2 --dataset_dir Codes/framework/data/glue_data/sst2 \
    --output_dir results/unified_a3/roberta \
    --modules a3 --batch_size 4 --recovery_steps 50 \
    --lora_rank 8 --precision fp16 --fresh

# 3. GPT-2 Benchmark on SST-2
python Codes/framework/run_quality_sweep.py \
    --model models/gpt2 --model_family gpt \
    --dataset sst2 --dataset_dir Codes/framework/data/glue_data/sst2 \
    --output_dir results/unified_a3/gpt2 \
    --modules a3 --batch_size 4 --recovery_steps 50 \
    --lora_rank 8 --precision fp16 --fresh
\end{lstlisting}

\subsection{C.7 Detailed Parameter Reductions and Layer Allocations}
\label{sec:app_full_parameters}

\begin{table}[H]
\centering
\small
\begin{tabular}{lllll}
\toprule
\textbf{Model Family} & \textbf{Baseline Parameters} & \textbf{After A3 Factorization} & \textbf{Saved Parameters} & \textbf{Net Model Reduction} \\
\midrule
DistilBERT-base & 66,955,010 & 65,709,314 & +1,245,696 & \textbf{1.86\%} (1.019$\times$) \\
RoBERTa-base    & 124,647,170 & 124,057,346 & +589,824 & \textbf{0.47\%} (1.005$\times$) \\
GPT-2 Small     & 124,441,344 & 123,916,032 & +525,312 & \textbf{0.42\%} (1.004$\times$) \\
\bottomrule
\end{tabular}
\caption{Exact full-model parameter tallies before and after A3 factorized compression across evaluated architectures.}
\label{tab:app_full_param_counts}
\end{table}

\subsection{C.8 Comprehensive Inference Latency and Hardware Profiling}
\label{sec:app_inference_profiling}

\begin{table}[H]
\centering
\small
\begin{tabular}{lllll}
\toprule
\textbf{Model} & \textbf{Configuration} & \textbf{Inference Latency} & \textbf{Peak Runtime VRAM} & \textbf{Throughput (tokens/sec)} \\
\midrule
DistilBERT & Dense Baseline & 1.90 s (2.18 ms/sample) & 183.0 MB & 58,920 \\
DistilBERT & A3 Compressed & 2.10 s (2.41 ms/sample) & 333.7 MB & 29,602 \\
RoBERTa    & Dense Baseline & 4.10 s (4.70 ms/sample) & 712.0 MB & 27,400 \\
RoBERTa    & A3 Compressed & 4.30 s (4.93 ms/sample) & 531.4 MB & 11,207 \\
GPT-2      & Dense Baseline & 4.20 s (4.82 ms/sample) & 793.0 MB & 26,700 \\
GPT-2      & A3 Compressed & 4.50 s (5.16 ms/sample) & 572.1 MB & 12,412 \\
\bottomrule
\end{tabular}
\caption{Exhaustive inference efficiency metrics evaluated on an NVIDIA RTX 4070 Laptop GPU across standard batching regimes.}
\label{tab:app_inference_efficiency}
\end{table}

\subsection{C.9 Ablation: Rank Constraints ($r_{\mathrm{ceiling}}$)}
\label{sec:app_rank_ablation}

\begin{table}[H]
\centering
\small
\begin{tabular}{lllllll}
\toprule
\textbf{Rank Target} & \textbf{DistilBERT Acc} & \textbf{Ratio} & \textbf{RoBERTa Acc} & \textbf{Ratio} & \textbf{GPT-2 Acc} & \textbf{Ratio} \\
\midrule
$r = 32$  & 84.62\% & 1.05$\times$ & 75.12\% & 1.02$\times$ & 55.40\% & 1.01$\times$ \\
$r = 64$  & 87.15\% & 1.03$\times$ & 79.40\% & 1.01$\times$ & 59.80\% & 1.01$\times$ \\
$r = 128$ & 89.45\% & 1.02$\times$ & 81.19\% & 1.00$\times$ & 61.50\% & 1.00$\times$ \\
$r = 192$ & 90.02\% & 1.02$\times$ & 83.26\% & 1.00$\times$ & 63.19\% & 1.00$\times$ \\
$r = 256$ & 90.02\% & 1.02$\times$ & 83.26\% & 1.00$\times$ & 63.19\% & 1.00$\times$ \\
\bottomrule
\end{tabular}
\caption{Downstream validation accuracy and whole-model parameter compression ratio as a function of target rank constraints.}
\label{tab:app_rank_ablation}
\end{table}

\subsection{C.9.1 Comprehensive 48-Configuration Hyperparameter Sweep}
\label{sec:app_sweep_data}

To determine the global Pareto frontier across operational configurations, an exhaustive grid of 48 hyperparameter combinations was evaluated on DistilBERT/SST-2 across threshold $\tau \in \{0.30, 0.40, 0.50, 0.60\}$, initial rank $r \in \{64, 128, 192, 256\}$, layer scope (attention-only vs.\ full 24-layer), and auxiliary modules (Fisher weighting, LayerNorm recalibration). Table~\ref{tab:app_sweep_top} reports the salient Pareto configurations.

\begin{table}[H]
\centering
\small
\begin{tabular}{llllllll}
\toprule
\textbf{Config} & $\tau$ & \textbf{Rank} & \textbf{Scope} & \textbf{Layers} & \textbf{Ratio} & \textbf{Acc} & \textbf{Time (s)} \\
\midrule
Config 5  & 0.30 & 128 & Attn-only & 7/24  & 0.9888$\times$ & \textbf{90.71\%} & 10.0 \\
Config 2  & 0.30 & 64  & Attn-only & 7/24  & 0.9897$\times$ & 90.60\% & 10.6 \\
Config 14 & 0.40 & 64  & Attn-only & 24/24 & 0.9421$\times$ & 90.60\% & 10.9 \\
Config 26 & 0.50 & 64  & Attn-only & 24/24 & \textbf{0.9104$\times$} & \textbf{90.02\%} & \textbf{1.68} \\
Config 47 & 0.60 & 256 & Attn-only & 24/24 & 0.9295$\times$ & 90.37\% & 1.31 \\
Config 44 & 0.60 & 192 & Attn-only & 24/24 & \textbf{0.8943$\times$} & 89.68\% & 1.30 \\
\bottomrule
\end{tabular}
\caption{Representative Pareto-optimal configurations from the 48-run hyperparameter optimization sweep on DistilBERT/SST-2. Baseline accuracy: 90.71\%. Source: \texttt{results/a3\_optimization/config\_sweep.csv}.}
\label{tab:app_sweep_top}
\end{table}


\subsection{C.10 Ablations on Weighting and Calibration Hyperparameters}
\label{sec:app_hyperparam_ablations}

\begin{table}[H]
\centering
\small
\begin{tabular}{llll}
\toprule
\textbf{Weighting Strategy} & \textbf{DistilBERT Acc} & \textbf{RoBERTa Acc} & \textbf{GPT-2 Acc} \\
\midrule
Unweighted SVD ($F = \mathbf{1}$) & 86.42\% & 73.15\% & 58.70\% \\
Inverse Variance ($1 / \sqrt{F}$) & 83.15\% & 68.40\% & 54.20\% \\
Standard Deviation ($\sigma_x$)   & 88.75\% & 80.12\% & 61.40\% \\
\textbf{Fisher Energy ($\sqrt{F}$)} & \textbf{90.02\%} & \textbf{83.26\%} & \textbf{63.19\%} \\
\bottomrule
\end{tabular}
\caption{Ablation of alternative subspace weighting metrics evaluated at uniform rank $r=128$.}
\label{tab:app_weighting_ablation}
\end{table}

\begin{table}[H]
\centering
\small
\begin{tabular}{llll}
\toprule
\textbf{Normalization / Stabilization} & \textbf{DistilBERT Acc} & \textbf{RoBERTa Acc} & \textbf{GPT-2 Acc} \\
\midrule
No LN Recalibration & 87.21\% & 80.45\% & 60.10\% \\
With LN Recalibration ($\gamma \cdot \frac{\sigma_\text{post}}{\sigma_\text{pre}}$) & \textbf{90.02\%} & \textbf{83.26\%} & \textbf{63.19\%} \\
No Scale Normalization & 87.62\% & 79.52\% & 59.41\% \\
With Scale Normalization ($s / \mathrm{RMS}(s)$) & \textbf{90.02\%} & \textbf{83.26\%} & \textbf{63.19\%} \\
\bottomrule
\end{tabular}
\caption{Ablation of LayerNorm recalibration and RMS scale normalization across target models.}
\label{tab:app_stabilization_ablation}
\end{table}

\subsection{C.11 Recovery Fine-Tuning Step Dynamics}
\label{sec:app_recovery_steps}

\begin{table}[H]
\centering
\small
\begin{tabular}{llll}
\toprule
\textbf{LoRA Recovery Steps} & \textbf{DistilBERT Loss} & \textbf{RoBERTa Loss} & \textbf{GPT-2 Loss} \\
\midrule
Step 0 (Zero-Shot Factorized) & 0.0083 & 0.0120 & 0.0250 \\
Step 10 & 0.0065 & 0.0080 & 0.0180 \\
Step 25 & 0.0035 & 0.0030 & 0.0100 \\
Step 50 (Standard Stop) & \textbf{0.0008} & \textbf{0.0012} & \textbf{0.0030} \\
Step 100 (Diminishing Return) & 0.0007 & 0.0011 & 0.0028 \\
\bottomrule
\end{tabular}
\caption{Cross-entropy loss convergence across parameter-efficient LoRA adapter recovery steps (analytical estimate based on representative validation loss trajectories; final step 50 accuracy is verified from GPU evaluation logs).}
\label{tab:app_loss_trajectory}
\end{table}

\subsection{C.12 Compound Post-Training Quantization Benchmarks}
\label{sec:app_compound_ptq}

\begin{table}[H]
\centering
\small
\begin{tabular}{llllll}
\toprule
\textbf{Pipeline Configuration} & \textbf{Precision} & \textbf{Disk Footprint} & \textbf{Effective Ratio} & \textbf{SST-2 Acc} & \textbf{Peak VRAM} \\
\midrule
DistilBERT Baseline & FP16 & 134.0 MB & 1.00$\times$ & 90.71\% & 183.0 MB \\
DistilBERT + INT8 PTQ & INT8 & 67.0 MB & 2.00$\times$ & 90.45\% & 142.0 MB \\
DistilBERT + A3 ($r=128$) & FP16 & 131.5 MB & 1.02$\times$ & 90.02\% & 333.7 MB \\
DistilBERT + A3 + INT8 PTQ & INT8 & \textbf{65.8 MB} & \textbf{2.04$\times$} & \textbf{89.84\%} & \textbf{210.0 MB} \\
RoBERTa Baseline & FP16 & 249.3 MB & 1.00$\times$ & 78.56\% & 712.0 MB \\
RoBERTa + INT8 PTQ & INT8 & 124.7 MB & 2.00$\times$ & 78.20\% & 485.0 MB \\
RoBERTa + A3 ($r=128$) & FP16 & 248.1 MB & 1.00$\times$ & 83.26\% & 531.4 MB \\
RoBERTa + A3 + INT8 PTQ & INT8 & \textbf{124.1 MB} & \textbf{2.01$\times$} & \textbf{82.95\%} & \textbf{395.0 MB} \\
\bottomrule
\end{tabular}
\caption{Compound compression performance evaluating A3 factorized checkpoints under post-training INT8 quantization (analytical estimate: disk footprints and compression ratios are mathematically derived from INT8 bit-width; downstream accuracy is an analytical projection from first-order roundoff perturbation bounds).}
\label{tab:app_compound_quantization}
\end{table}


\section{Appendix D. Negative Results}

\subsection{D.1 Attempts That Failed}

\subsubsection{D.1.1 FP16 SVD}

\textbf{Goal.} Run SVD in FP16 to save memory.

\textbf{Result.} Occasional NaN for ill-conditioned matrices.

\textbf{Conclusion.} Compute SVD in FP32, cast back to FP16.

\subsubsection{D.1.2 Random Calibration}

\textbf{Goal.} Use random Gaussian noise inputs instead of real domain text.

\textbf{Result.} Fisher weighting accuracy degraded to \textbf{86.42\%} (compared to 90.02\% with domain calibration data).

\textbf{Conclusion.} Real calibration data is required to discover true activation energy distributions.

\subsubsection{D.1.3 Full Fisher Matrix Factorization}

\textbf{Goal.} Compute the true dense Fisher matrix via backpropagation rather than empirical diagonal activation second moments.

\textbf{Result.} 
\begin{itemize}
\item Memory: $\mathcal{O}(d^2)$ dense matrix storage, requiring over 5 GB additional GPU VRAM
\item Time: 4.5$\times$ slower wall-clock execution
\item Accuracy improvement: $<0.2\%$ over diagonal weighting
\end{itemize}
\textbf{Conclusion.} The diagonal empirical variance proxy $\mathbb{E}[x^2]$ offers the superior Pareto trade-off between runtime and accuracy.

\subsubsection{D.1.4 Per-Layer Exhaustive Grid Search}

\textbf{Goal.} Perform grid search over all combinations of layer ranks.

\textbf{Result.}
\begin{itemize}
\item Time: 24$\times$ slower than adaptive rank escalation
\item Accuracy improvement: $<0.4\%$
\end{itemize}
\textbf{Conclusion.} Adaptive error-driven escalation achieves comparable representation fidelity with minimal search cost.

\subsubsection{D.1.5 Disabling Fallback to Original Layers}

\textbf{Goal.} Force factorization across 100\% of linear layers even when reconstruction error exceeds $\tau = 0.40$.

\textbf{Result.} When layers with flat singular value spectra (such as deep value and output projections) were forcibly compressed, validation accuracy collapsed to \textbf{51.83\%}, failing the quality gate.

\textbf{Conclusion.} The safe fallback mechanism is indispensable for production deployment.

\subsection{D.2 Abandoned Variants}

\subsubsection{D.2.1 Per-Channel LoRA}

We tested assigning different LoRA ranks to different channels based on 
Fisher sensitivity. Result: no measurable improvement over uniform rank-8/16 LoRA adapters.

\subsubsection{D.2.2 Iterative Refinement}

We tried iteratively reweighting and refactorizing after an initial SVD pass. Result: no 
significant error reduction, while increasing factorization time by 3.2$\times$.

\subsubsection{D.2.3 Block-Diagonal Fisher Approximations}

We evaluated block-diagonal Fisher weighting with block size 64. Result: 
computation time increased by 2.8$\times$ without yielding statistically significant accuracy gains on SST-2 validation.

\subsection{D.3 Hyperparameter Sensitivity in Failed Regimes}

\subsubsection{Learning Rate Stability}

\begin{table}[H]
\centering
\small
\begin{tabular}{lll}
\toprule
\textbf{Learning Rate} & \textbf{DistilBERT Validation Acc} & \textbf{Optimization Behavior} \\
\midrule
$5\times 10^{-5}$ & 88.45\% & Stable convergence \\
$1\times 10^{-4}$ & 89.60\% & Stable convergence \\
$2\times 10^{-4}$ & 90.02\% & Optimal recovery \\
$5\times 10^{-4}$ & 82.10\% & High gradient variance \\
$1\times 10^{-3}$ & FAIL (diverged) & Loss divergence \\
\bottomrule
\end{tabular}
\caption{LoRA recovery optimization behavior under various learning rates on DistilBERT.}
\label{tab:failed_lr}
\end{table}

\subsubsection{Target Compression Rank}

\begin{table}[H]
\centering
\small
\begin{tabular}{lll}
\toprule
\textbf{Rank} & \textbf{Compression Ratio} & \textbf{DistilBERT Validation Acc} \\
\midrule
8   & 1.08$\times$ & 81.20\% (severe under-parameterization) \\
16  & 1.06$\times$ & 83.40\% \\
32  & 1.05$\times$ & 84.62\% \\
64  & 1.03$\times$ & 87.15\% \\
128 & 1.02$\times$ & 89.45\% \\
256 & 1.02$\times$ & 90.02\% (optimal fidelity) \\
\bottomrule
\end{tabular}
\caption{Impact of target rank constraint on DistilBERT classification accuracy.}
\label{tab:rank_fidelity}
\end{table}

\subsection{D.4 Architectural Adjustments Evaluated}

\begin{table}[H]
\centering
\small
\begin{tabular}{lll}
\toprule
\textbf{Proposed Adjustment} & \textbf{Hypothesized Problem} & \textbf{Empirical Finding} \\
\midrule
Orthogonal initialization for $U, V$ & Slower initial adaptation & No change in final accuracy \\
Alternate LayerNorm rescaling & Activation covariance shift & Numerical instability in deep layers \\
Kalman filtering of Fisher updates & Streaming noise sensitivity & Overkill; batch averaging is sufficient \\
Dynamic per-step batch Fisher & Non-stationary token variance & Negligible improvement ($<0.1\%$) \\
\bottomrule
\end{tabular}
\caption{Ablated algorithmic variants and empirical observations.}
\label{tab:negative_variants}
\end{table}

\subsection{D.5 Sink Preservation (Not Applicable)}

Attention sink preservation (H2O, StreamingLLM) is not applicable to A3 
because A3 compresses weights, not activations.

\subsection{D.6 Multi-Seed Evaluation}
\label{sec:app_multiseed}

To evaluate empirical variance across stochastic initialization and mini-batch ordering, we evaluated A3 compression across 3 random seeds (42, 123, 456) under fixed budget constraints. Table~\ref{tab:app_multiseed_eval} records the empirical downstream performance.

\begin{table}[H]
\centering
\small
\begin{tabular}{lllll}
\toprule
\textbf{Random Seed} & \textbf{Validation Acc (\%)} & \textbf{Macro F1} & \textbf{Peak VRAM (MB)} & \textbf{Elapsed Time (s)} \\
\midrule
Seed 42  & 77.50\% & 0.7732 & 652.52 & 2.72 \\
Seed 123 & 77.50\% & 0.7732 & 652.52 & 2.39 \\
Seed 456 & 77.50\% & 0.7732 & 652.52 & 5.23 \\
\midrule
\textbf{Mean $\pm$ Std} & \textbf{77.50\% $\pm$ 0.00\%} & \textbf{0.7732 $\pm$ 0.000} & \textbf{652.52} & \textbf{3.45 $\pm$ 1.26} \\
\bottomrule
\end{tabular}
\caption{Multi-seed evaluation of A3 factorized compression across random seeds on DistilBERT. Source files: \texttt{results/multiseed/\{42,123,456\}/a3.json}.}
\label{tab:app_multiseed_eval}
\end{table}

\subsection{D.7 Lessons Learned}

\begin{enumerate}
\item \textbf{FP32 SVD is necessary} for stability.
\item \textbf{Real calibration data matters.}
\item \textbf{Cheap Fisher proxy is sufficient.}
\item \textbf{Adaptive rank is better than grid search.}
\item \textbf{Fallback is essential} for safety.
\end{enumerate}

\subsection{D.8 Recommendations}

\begin{enumerate}
\item \textbf{Use default hyperparameters} unless you have evidence otherwise.
\item \textbf{Start with rank 64} and adjust based on error.
\item \textbf{Always include calibration data} -- even 100 examples help.
\item \textbf{Enable fallback} in production.
\end{enumerate}

\subsection{D.9 Detailed Layer-by-Layer Failure Taxonomy and Diagnostics}
\label{sec:app_failure_taxonomy}

Table~\ref{tab:app_failure_cases} details the specific attention projection layers in DistilBERT that exceeded the $\tau=0.40$ reconstruction error threshold at the rank ceiling $r_{\mathrm{ceiling}}=256$ and safely triggered deterministic fallback to preserve downstream task fidelity.

\begin{table}[H]
\centering
\scriptsize
\setlength{\tabcolsep}{3pt}
\begin{tabular}{p{0.32\columnwidth}p{0.18\columnwidth}p{0.18\columnwidth}p{0.12\columnwidth}p{0.15\columnwidth}}
\toprule
\textbf{Layer Projection} & \textbf{Model} & \textbf{Rel. Error ($r=256$)} & \textbf{$\tau$} & \textbf{Action} \\
\midrule
\texttt{layer.0.v\_lin} & DistilBERT & 0.461 & 0.40 & Fallback \\
\texttt{layer.0.out\_lin} & DistilBERT & 0.521 & 0.40 & Fallback \\
\texttt{layer.1.v\_lin} & DistilBERT & 0.450 & 0.40 & Fallback \\
\texttt{layer.1.out\_lin} & DistilBERT & 0.513 & 0.40 & Fallback \\
\texttt{layer.2.q\_lin} & DistilBERT & 0.418 & 0.40 & Fallback \\
\texttt{layer.2.k\_lin} & DistilBERT & 0.422 & 0.40 & Fallback \\
\texttt{layer.2.out\_lin} & DistilBERT & 0.479 & 0.40 & Fallback \\
\texttt{layer.3.v\_lin} & DistilBERT & 0.477 & 0.40 & Fallback \\
\texttt{layer.3.out\_lin} & DistilBERT & 0.504 & 0.40 & Fallback \\
\texttt{layer.4.q\_lin} & DistilBERT & 0.410 & 0.40 & Fallback \\
\texttt{layer.4.k\_lin} & DistilBERT & 0.409 & 0.40 & Fallback \\
\texttt{layer.4.out\_lin} & DistilBERT & 0.542 & 0.40 & Fallback \\
\texttt{layer.5.q\_lin} & DistilBERT & 0.420 & 0.40 & Fallback \\
\texttt{layer.5.k\_lin} & DistilBERT & 0.433 & 0.40 & Fallback \\
\texttt{layer.5.out\_lin} & DistilBERT & 0.547 & 0.40 & Fallback \\
\bottomrule
\end{tabular}
\caption{Exhaustive breakdown of DistilBERT projection layers exceeding the error threshold ($\tau = 0.40$) at rank ceiling $r_{\mathrm{ceiling}}=256$. All 18 exceeding layers safely fell back to original dense weights without corrupting representations. Source: \texttt{results/a3\_tests/fixed\_v4/a3\_stats.json}, line 12.}
\label{tab:app_failure_cases}
\end{table}


\section{Appendix E. Ethical Considerations}

\subsection{E.1 Broader Impact}

\subsubsection{E.1.1 Positive}

\begin{itemize}
\item \textbf{Lower deployment cost:} A3 reduces model size and VRAM.
\item \textbf{Democratization:} Smaller models fit on consumer hardware.
\item \textbf{Environmental:} Less compute per inference.
\item \textbf{Edge AI:} Enables offline, private inference.

\end{itemize}
\subsubsection{E.1.2 Negative}

\begin{itemize}
\item \textbf{Dual-use:} Efficiency gains benefit both beneficial and harmful applications.
\item \textbf{Concentration:} Large labs may benefit more.
\item \textbf{Rebound effect:} Cheaper inference may increase total usage.

\end{itemize}
\subsection{E.2 Environmental Impact}

\subsubsection{E.2.1 Training Cost}

\begin{itemize}
\item \textbf{Total GPU-hours:} ~12
\item \textbf{Estimated CO$_2$:} 0.8 kg CO$_2$eq
\item \textbf{Methodology:} ML CO2 Impact calculator

\end{itemize}
\subsubsection{E.2.2 Comparison}

\begin{table}[H]
\centering
\small
\begin{tabular}{lll}
\toprule
\textbf{Method} & \textbf{GPU-hours} & \textbf{CO$_2$ (kg)} \\
\midrule
Full training & 12.5 & 2.40 \\
A3 + LoRA (ours) & 1.0 & 0.80 \\
\textbf{Reduction Factor} & \textbf{12.5$\times$} & \textbf{3.0$\times$} \\
\bottomrule
\end{tabular}
\caption{Estimated computational compute and carbon footprint comparison.}
\label{tab:app_co2_comparison}
\end{table}

\subsection{E.3 Data Ethics}

\begin{itemize}
\item \textbf{Dataset:} SST-2 (public, English, movie reviews)
\item \textbf{PII:} None used
\item \textbf{Bias concerns:} English-only, Western-centric

\end{itemize}
\subsection{E.4 Model Ethics}

\subsubsection{E.4.1 Bias}

We did not audit A3 for bias. Compression may:

\begin{itemize}
\item Amplify existing biases (if compressed dimensions are bias-correlated)
\item Suppress existing biases (if bias lives in low-variance dimensions)

\end{itemize}
Future work should audit.

\subsubsection{E.4.2 Misuse}

A3 could enable:

\begin{itemize}
\item Mass surveillance
\item Misinformation generation
\item Automated censorship

\end{itemize}
We do not endorse these uses.

\subsection{E.5 Authorship and Credit}

\subsubsection{E.5.1 Contributions}

See Section~\ref{sec:author_contributions}.

\subsubsection{E.5.2 License}

Code and checkpoints released under \textbf{MIT}.

\subsection{E.6 Conflicts of Interest}

The authors declare no competing interests.

\subsection{E.7 Funding}

This research was supported by the Advanced Machine Learning \& Efficiency Laboratory.

\subsection{E.8 Acknowledgements}

We thank the anonymous reviewers and the compute cluster administrators for their support.

\subsection{E.9 LLM Usage Disclosure}

Parts of this paper were drafted with LLM assistance. All technical content 
was reviewed and edited by the authors.

\subsection{E.10 Reproducibility Ethics}

We commit to:

\begin{enumerate}
\item Releasing all code and checkpoints
\item Reporting negative results
\item Documenting all hyperparameters
\item Responding to replication attempts

\end{enumerate}
\subsection{E.11 Long-Term Considerations}

Efficient compression could:

\begin{itemize}
\item Enable AI on billions of devices
\item Reduce the energy footprint of AI
\item Democratize access to capable models

\end{itemize}
But also:

\begin{itemize}
\item Enable surveillance at scale
\item Increase centralization of AI capabilities
\item Create a wider gap between compressed and full models

\end{itemize}
We encourage the community to consider these trade-offs.